% The document class supplies options to control rendering of some standard
% features in the result.  The goal is for uniform style, so some attention 
% to detail is *vital* with all fields.  Each field (i.e., text inside the
% curly braces below, so the MEng text inside {MEng} for instance) should 
% take into account the following:
%
% - author name       should be formatted as "FirstName LastName"
%   (not "Initial LastName" for example),
% - supervisor name   should be formatted as "Title FirstName LastName"
%   (where Title is "Dr." or "Prof." for example),
% - degree programme  should be "BSc", "MEng", "MSci", "MSc" or "PhD",
% - dissertation title should be correctly capitalised (plus you can have
%   an optional sub-title if appropriate, or leave this field blank),
% - dissertation type should be formatted as one of the following:
%   * for the MEng degree programme either "enterprise" or "research" to
%     reflect the stream,
%   * for the MSc  degree programme "$X/Y/Z$" for a project deemed to be
%     X%, Y% and Z% of type I, II and III.
% - year              should be formatted as a 4-digit year of submission
%   (so 2014 rather than the academic year, say 2013/14 say).

\documentclass[
    author={Aanchala Ravindra Bhongade},
    supervisor={Dr. Fernando Moreno-Pino},
    degree={MSc},
    title={Length Generalisation in Transformer-Based Sequence Models: Failure Taxonomy and Mitigation Strategies},
    subtitle={},
    type={Research},
    year={2026}
]{dissertation}

\renewcommand{\today}{Friday 4 September, 2026}

\graphicspath{{figures/}}

\begin{document}


\maketitle
\frontmatter
\makedecl

% ----------------------------------------------------------------------------- ACKNOLEDGEMENT---------

\chapter*{Acknowledgements}

\noindent
I would like to express my sincere gratitude to my supervisor, Dr. Fernando Moreno-Pino,  for his guidance, feedback, and support throughout this project. His advice helped me refine the experimental design, interpret the results more carefully, and develop a clearer direction for the dissertation. I am equally grateful to my personal tutor, Dr. Conor Houghton, for his support and encouragement throughout my journey at the University of Bristol.

I would also like to acknowledge the University of Bristol for providing the resources and academic environment that made this work possible.

Finally, I am grateful to my family and friends for their encouragement and support during the completion of this dissertation.




%==================================

\tableofcontents
\listoffigures
\listoftables
%\listofalgorithms
%\lstlistoflistings

% The following sections are part of the front matter, but are not generated
% automatically by LaTeX; the use of \chapter* means they are not numbered.



% -----------------------------------------------------------------------------

\chapter*{Supporting Technologies}
\addcontentsline{toc}{chapter}{Supporting Technologies}

\begin{itemize}
    \item Python was used to implement dataset generation, model training,
    evaluation, diagnostic analysis, and result visualisation.

    \item PyTorch was used to implement the Transformer models, positional
    encoding variants, optimisation procedure, checkpointing, and GPU
    execution.

    \item NumPy and pandas were used for numerical processing, aggregation
    of experimental results, and construction of reproducible summary tables.

    \item Matplotlib was used to produce the dissertation figures. A common
    visual style, colour mapping, typography, and output format were applied
    across the final plots.

    \item PyYAML configuration files were used to define experiment settings
    separately from the training and evaluation implementation.

    \item The aeon dataset tools were used to obtain and prepare the publicly
    available SelfRegulationSCP2 time-series dataset.

    \item Amazon Web Services EC2 G5-family GPU resources were used for the
    multi-seed training and evaluation runs that were impractical to execute
    locally.

    \item Git and GitHub were used for version control and for synchronising
    the experimental implementation between local and AWS environments.

    \item LaTeX was used to prepare the dissertation and to include
    publication-quality PDF figures.
\end{itemize}

% -----------------------------------------------------------------------------

\chapter*{Notation and Acronyms}
\addcontentsline{toc}{chapter}{Notation and Acronyms}

\section*{Acronyms}

\begin{description}
    \item[ALiBi] Attention with Linear Biases.
    \item[AWS] Amazon Web Services.
    \item[EM] Exact match.
    \item[FFN] Feed-forward network.
    \item[GPU] Graphics Processing Unit.
    \item[ID] In-distribution evaluation.
    \item[LSTM] Long Short-Term Memory.
    \item[NoPE] No explicit positional encoding.
    \item[OOD] Out-of-distribution evaluation.
    \item[PE] Positional encoding.
    \item[RNN] Recurrent neural network.
    \item[RoPE] Rotary Position Embedding.
    \item[SD] Standard deviation.
    \item[UEA] University of East Anglia.
\end{description}

\section*{Mathematical Notation}

\begin{description}
    \item[$n$] Sequence length.

    \item[$d_{\mathrm{model}}$] Dimensionality of the Transformer
    representations.

    \item[$d_k$] Dimensionality of the key and query vectors within an
    attention head.

    \item[$N_{\mathrm{head}}$] Number of attention heads.

    \item[$\mathbf{X}$] Matrix of input representations supplied to a
    self-attention layer.

    \item[$\mathbf{Q},\mathbf{K},\mathbf{V}$] Query, key, and value matrices
    used in self-attention.

    \item[$\mathcal{L}_{\mathrm{train}}$] Set of sequence lengths represented
    during training.

    \item[$L_{\mathrm{train}}$] Sequence length used to train a fixed-length
    configuration.

    \item[$L_{\mathrm{test}}$] Sequence length used during evaluation.

    \item[$L_{\max}$] Maximum sequence length represented during training,
    $\max\mathcal{L}_{\mathrm{train}}$.

    \item[$A_{\mathrm{EM}}(L)$] Exact-match accuracy at test length $L$.

    \item[$A_{\mathrm{tok}}(L)$] Token accuracy at test length $L$.

    \item[$E_{\mathrm{EM}}(L)$] Exact-match error percentage,
    $100-A_{\mathrm{EM}}(L)$.

    \item[$E_{\mathrm{tok}}(L)$] Token-error percentage,
    $100-A_{\mathrm{tok}}(L)$.

    \item[$\mathcal{V}_{\mathrm{target}}$] Set of possible target-output
    symbols.

    \item[$|\mathcal{V}_{\mathrm{target}}|$] Number of possible target-output
    symbols.

    \item[$E_{\mathrm{chance}}$] Uniform-prediction token-error reference,
    $100\left(1-1/|\mathcal{V}_{\mathrm{target}}|\right)$.
\end{description}

% -----------------------------------------------------------------------------
\chapter*{Reproducibility Statement}
\addcontentsline{toc}{chapter}{Reproducibility Statement}

\noindent
The experimental implementation separates task generation, model
definitions, training, evaluation, configuration files, raw results, and
derived plots. Synthetic experiments were repeated using random seeds 42,
123, and 2024. Configuration files record the model architecture, positional
encoding, training length, evaluation lengths, optimiser settings, and output
locations. Final figures were generated from saved CSV results rather than
manually entered values. Representative position-wise diagnostics use the
seed-42 checkpoints and are explicitly distinguished from the multi-seed
baseline analysis.

% ----------------------------------------------------------------------------- ABSTRACT------------

\chapter*{Abstract}

\noindent
Transformer models have become a dominant architecture for sequence modelling, partly because self-attention allows them to represent dependencies across an input without the recurrence used in earlier sequence models. However, strong performance on the sequence lengths observed during training does not necessarily imply that the same model has learned a rule that generalises to longer sequences. This limitation is important because many practical sequence-modelling problems require models to operate reliably outside the exact length distribution seen during training.

This dissertation investigates length generalisation in Transformer-based sequence models using a controlled train-short/test-long experimental framework. The project focuses on three synthetic sequence tasks: addition, delayed copy, and reverse. These tasks were selected because they test different forms of sequence reasoning, including symbolic computation, long-range memory, and position-dependent transformations. A compact Transformer model was evaluated across five positional encoding settings: learned positional embeddings, sinusoidal positional encodings, no explicit positional encoding, ALiBi, and RoPE. Experiments were repeated across three random seeds to assess robustness rather than relying on single-run results.


A baseline-eligibility procedure first identified configurations that reliably
learned the training task. Eligible configurations were then analysed using
exact-match accuracy, token-error percentages, position-wise error profiles,
attention diagnostics, and task-specific tests of carry conditions and
dependency distance. Mixed-length training, curriculum learning, and
randomised padding were evaluated against matched original baselines to test
specific hypotheses about exposure to varied lengths and dependence on
absolute positions.

All five eligible task-encoding configurations that met the multi-seed baseline-eligibility criterion fell below the 10\% exact-match
threshold at the first unseen length. At this point, Addition showed
approximately 86.9--88.0\% token error, with errors increasing modestly
towards later output positions. Delayed Copy and Reverse showed approximately
98.8--99.0\% token error that was almost uniform across output positions.
Addition errors differed by less than one percentage point between digits
with and without an incoming carry for the learned model and were effectively
identical for the sinusoidal model. Reverse errors were not concentrated at
a particular dependency distance. Mixed-length variants underfit in the
evaluated settings, while curriculum and randomised-padding variants that
retained training-length performance produced no observed improvement at
unseen lengths. In the limited SelfRegulationSCP2 extension, mean
classification accuracy remained approximately 49--56\% across the evaluated
input lengths rather than showing the catastrophic exact-match collapse
observed on the synthetic tasks.

The resulting failure taxonomy distinguishes training-length underfitting,
seed instability, immediate extrapolation collapse, position-dependent error structure,
and no observed improvement under tested mitigation. These findings show
that length-generalisation evaluation should establish a valid training
baseline before examining not only whether performance decreases, but how
failure is distributed and which intervention hypothesis was actually tested.

\vfill

\noindent
\textbf{Keywords:} Length generalisation; Transformers; positional encoding;
synthetic sequence tasks; extrapolation; failure taxonomy; baseline
eligibility; position-wise diagnostics; mitigation strategies; time-series
classification.


% =============================================================================

% After the front matter comes a number of chapters; under each chapter,
% sections, subsections and even subsubsections are permissible.  The
% pages in this part are numbered with Arabic numerals.  Note that:
%
% - A reference point can be marked using \label{XXX}, and then later
%   referred to via \ref{XXX}; for example Chapter\ref{chap:context}.
% - The chapters are presented here in one file; this can become hard
%   to manage.  An alternative is to save the content in seprate files
%   the use \input{XXX} to import it, which acts like the #include
%   directive in C.

\mainmatter

% -----------------------------------------------------------------------------

\chapter{Introduction}
\label{chap:introduction}

\section{Context and Motivation}

\noindent
Sequence models are used in applications where the amount of available
context can vary substantially between examples. A model may be trained using
sequences of a particular length because of data availability, memory
constraints, or computational cost, but later be required to process longer
inputs. Reliable performance therefore requires more than fitting the length
distribution represented in the training data: the model must also transfer
the learned computation to sequence lengths that were not observed during
training.

Transformers have become a widely used architecture for sequence modelling
because self-attention allows direct interaction between positions without
the recurrent computation used in earlier sequence models
\cite{vaswani2017attention}. However, self-attention does not by itself
guarantee that a model will learn a length-independent rule. A Transformer can
achieve high accuracy within the training distribution while relying on
patterns that are specific to the positions or sequence lengths encountered
during training. Consequently, strong in-distribution performance may conceal
a substantial failure when the same model is evaluated on longer sequences.

Length generalisation refers to the ability of a model trained on shorter
instances to perform correctly on longer instances of the same underlying
task. It is a form of out-of-distribution generalisation because the task rule
remains unchanged while the sequence length moves beyond the training
distribution. Previous studies have shown that this remains difficult for
Transformer models, including on tasks for which a systematic algorithm
should, in principle, apply at arbitrary lengths
\cite{anil2022exploring,zhou2024algorithms}. This distinction is especially
important for arithmetic and algorithmic tasks: success on short examples
does not necessarily demonstrate that the model has learned the general
procedure needed to solve longer examples.

Positional encoding is closely connected to this problem because a
Transformer must represent token order while operating at positions that may
not have appeared during training. Learned absolute embeddings, fixed
sinusoidal encodings, rotary position embedding (RoPE) \cite{su2021roformer}, ALiBi \cite{press2022train}, and models without
explicit positional encodings (NoPE) \cite{kazemnejad2023impact} make different assumptions about how position
should be represented. Randomised positional encodings have specifically been proposed to expose models during training to positions beyond the ordinary range \cite{ruoss2023randomized}. Nevertheless,
published results vary across architectures, tasks, training procedures, and
evaluation metrics. A positional encoding that assists one experimental
setting cannot therefore be assumed to provide length generalisation
universally. This problem is related to, but distinct from, extending the usable context window or benchmarking efficient long-context architectures \cite{dai2019transformerxl,tay2021longrange}.

A further difficulty concerns how failure is evaluated. Reporting only
accuracy at the longest test length does not establish whether the model
learned the original training task, where extrapolation first deteriorated,
or how errors are distributed within the output. For example, a configuration
that never learns the training length represents training underfitting rather
than length-extrapolation failure. Similarly, two models with zero
sequence-level exact-match accuracy may exhibit different token-error
percentages and different positional error structures. These behaviours
require different interpretations even though their final exact-match scores
are identical.

This dissertation addresses these limitations through a controlled
train-short/test-long investigation of compact Transformer models. Addition,
delayed copy, and reverse tasks are used to represent symbolic computation,
long-range retention, and position-dependent transformation. The experiments
compare five positional encoding settings across three random seeds, first
establishing which configurations learn the training task reliably and then
examining how eligible configurations behave at progressively longer
sequence lengths. The study combines exact-match and token-level metrics with
position-wise, attention-based, and task-specific diagnostics. It also
evaluates targeted training interventions against matched original baselines
and includes a limited real-world time-series extension. The inclusion of this extension is motivated by the use of specialised Transformer architectures for long and irregularly sampled time-series data \cite{morenopino2024rough}.

\section{Research Problem, Aim, and Questions}
\label{sec:research_questions}

Existing studies demonstrate that Transformer length generalisation is
affected by positional representation, training procedure, task structure,
and output format \cite{anil2022exploring,kazemnejad2023impact,
ruoss2023randomized}. However, identifying which configuration achieves the
highest accuracy does not fully characterise the problem. A model may fail
because it never learned the training task, because its performance depends
strongly on random initialisation, or because a solution that works at the
training length does not transfer to longer sequences. These cases can
produce similarly poor long-length accuracy while representing different
experimental behaviours.

The central problem addressed in this dissertation is therefore not simply
whether Transformer performance decreases as sequence length increases. It is
how to distinguish the observable forms of failure, identify where errors
occur within the predicted output, and determine whether an intervention
changes the relevant behaviour when compared with an appropriate baseline.

In this dissertation, a failure taxonomy means a structured classification of empirically observed failure behaviours. The purpose of the taxonomy is not to assign a single label to every unsuccessful model, but to separate distinct experimental outcomes that have different interpretations, such as underfitting, extrapolation collapse, seed instability, and failure to improve under mitigation.

Importantly, negative results are treated as diagnostic evidence rather than as unsuccessful experiments. A poor result is interpreted according to whether the model fails to learn the training task, learns the training task but collapses beyond it, varies substantially across random seeds, or remains unchanged after a targeted intervention. \\

The primary research question is:

\begin{quote}
\itshape
When a Transformer model trained on short sequences is evaluated on
substantially longer sequences, which empirically distinguishable failure
behaviours occur, where within the output do the errors arise, and do targeted
training interventions alter these behaviours?
\end{quote}

This question is divided into five supporting questions:

\begin{enumerate}
    \item Which combinations of task and positional encoding reliably learn
    the training-length task across multiple random seeds?

    \item For configurations that learn the training task, at which unseen
    sequence length does performance first collapse, and how rapidly does it
    deteriorate thereafter?

    \item How are errors distributed across output positions, carry
    conditions, and dependency distances, and do these distributions differ
    between addition, delayed copy, and reverse?

    \item Do learned, sinusoidal, rotary, ALiBi, and no explicit positional
    encodings exhibit different training-fit, seed-stability, or
    extrapolation behaviours under otherwise matched conditions?

    \item Do mixed-length training, curriculum learning, and randomised
    padding improve unseen-length performance relative to matched original
    baselines while preserving training-length performance?
\end{enumerate}

The corresponding aim is to construct an evidence-based taxonomy that
separates training-length underfitting, seed instability, immediate
extrapolation collapse, position-dependent error structure, and no observed
improvement under tested mitigation. The purpose of this taxonomy is to
provide a more precise vocabulary for reporting length-generalisation
experiments than a single statement that a model ``fails'' at longer
sequences.

\section{Scope and Experimental Boundaries}
\label{sec:scope}

This dissertation investigates length generalisation using compact
Transformer models under controlled train-short/test-long conditions. The
main synthetic benchmark is restricted to addition, delayed copy, and
reverse, which provide deterministic tasks with controllable sequence
lengths. Associative recall was explored during development but excluded from
the main comparison because the model did not learn the training task
reliably enough to support a valid extrapolation analysis.

The study compares learned, sinusoidal, RoPE, ALiBi, and no explicit
positional encodings across three random seeds. Multi-seed experiments are
used to establish training fit, seed stability, and baseline eligibility.
More detailed position-wise and task-specific analyses use representative
seed-42 checkpoints and are explicitly distinguished from the multi-seed
results.

The mitigation study is limited to mixed-length training, curriculum
learning, and randomised padding. Each intervention is compared with a
matched original baseline and is evaluated according to whether it preserves
training-length performance while improving performance at unseen lengths.
The conclusions therefore concern the tested configurations rather than the
entire families of these training methods.

A limited real-world extension uses the SelfRegulationSCP2 time-series
classification dataset. Its inclusion is motivated by the broader use of
specialised Transformer architectures for long and irregularly sampled
time-series data \cite{morenopino2024rough}. However, this extension is
exploratory and is not presented as a state-of-the-art time-series benchmark
or as direct validation of the synthetic failure taxonomy.

The dissertation does not study parameter scaling, pre-trained language
models, retrieval systems, recurrent-memory architectures, or specialised
context-window extension methods. Attention and error diagnostics are used to
characterise observable behaviour, but they are not treated as complete
causal explanations of the models' internal computations.

\section{Contributions and Achievements}
\label{sec:contributions}

The contribution of this dissertation is primarily empirical and
methodological. It does not propose a new Transformer architecture. Instead,
it develops a controlled framework for distinguishing different forms of
poor long-length performance and applies that framework consistently across
tasks, positional encodings, random seeds, and training interventions.

The main contribution is therefore an evidence-based failure taxonomy supported by controlled experiments, multi-seed robustness analysis, position-wise diagnostics, attention diagnostics, and matched mitigation comparisons. \\

The principal contributions and achievements are:

\begin{itemize}
    \item A reproducible synthetic benchmark covering addition, delayed copy,
    and reverse with five positional encoding settings and three random
    seeds, producing 45 main synthetic training runs, alongside 15 real-world extension runs and separate mitigation experiments.

    \item A baseline-eligibility procedure that separates configurations
    which reliably learn the training task from training underfitting and
    seed-dependent success.

    \item An evaluation framework that combines exact-match accuracy with
    token-error percentages, degradation curves, and first-unseen-length
    comparisons.

    \item Position-wise and task-specific diagnostics that examine where
    errors occur, including output-position profiles, carry-conditioned
    addition errors, and dependency-distance errors for reverse.

    \item A matched comparison of mixed-length training, curriculum learning,
    and randomised padding against their corresponding original baselines,
    with each intervention interpreted according to its intended target.

    \item An evidence-based failure taxonomy that distinguishes
    training-length underfitting, seed instability, immediate extrapolation
    collapse, position-dependent error structure, and no observed improvement
    under tested mitigation.

    \item A limited real-world extension using SelfRegulationSCP2 to contrast
    synthetic sequence-generation failures with behaviour on a time-series
    classification task.

    \item A reproducible experimental implementation with separate configurations, checkpoints, raw CSV results, derived summaries, diagnostic outputs, and final report figures for traceability.
\end{itemize}

The main empirical outcome is that every task-encoding configuration meeting
the multi-seed baseline-eligibility criterion collapsed at the first unseen
length, although the token-error structure differed across tasks. The tested
mitigation variants either underfit or produced no observed improvement at
unseen lengths when training-length performance was preserved. These findings
motivate the dissertation's emphasis on describing the form and location of
failure rather than reporting only that accuracy decreases.

\section{Dissertation Structure}
\label{sec:dissertation_structure}

The remainder of this dissertation is organised as follows. Chapter 2
introduces the technical background and reviews existing research on
Transformers, positional encoding, length generalisation, algorithmic
evaluation, and long-context sequence modelling. Chapter 3 describes the
experimental methodology, including the synthetic tasks, model architecture,
baseline-eligibility procedure, evaluation metrics, diagnostic analyses,
mitigation experiments, and real-world extension.

Chapter 4 presents the multi-seed synthetic benchmark results and establishes
which task--encoding configurations provide valid extrapolation baselines.
Chapter 5 analyses the observed failures in greater detail using positional,
task-specific, and attention-based diagnostics before comparing the tested
mitigation strategies and constructing the failure taxonomy. Chapter 6
reports the limited SelfRegulationSCP2 time-series extension.

Chapter 7 discusses the findings in relation to the research questions and
existing literature, together with the study's limitations and implications.
Finally, Chapter 8 summarises the principal conclusions and identifies
directions for future work.
% -----------------------------------------------------------------------------

\chapter{Background}
\label{chap:background}

\section{Sequence Modelling and Transformer Architecture}
\label{sec:transformer_background}

Sequence modelling concerns data whose interpretation depends on the ordering
and interaction of its elements. Given an input sequence
$\mathbf{x}=(x_1,\ldots,x_n)$, a sequence model constructs representations that
can be used to predict an output at each position, classify the complete
sequence, or generate another sequence. Earlier neural sequence-to-sequence
systems commonly used recurrent architectures, in which a hidden state was
updated sequentially as each input element was processed
\cite{sutskever2014sequence}. This recurrence provided an explicit notion of
order, but it also created a sequential computational path through the input.

Attention mechanisms introduced a different way of representing dependencies.
Instead of relying exclusively on a single recurrent state, an attention
mechanism computes a weighted combination of representations according to
their relevance to the current prediction \cite{bahdanau2015neural}. The
Transformer extended this principle by replacing recurrence with stacked
self-attention and feed-forward layers \cite{vaswani2017attention}. As a
result, representations at different sequence positions can be updated in
parallel, while each position can directly incorporate information from other
positions.

\subsection{Sequence Representation and Self-Attention}
\label{subsec:self_attention}

Let the input to a self-attention layer be represented by
$\mathbf{X}\in\mathbb{R}^{n\times d_{\mathrm{model}}}$, where $n$ is the
sequence length and $d_{\mathrm{model}}$ is the representation dimension.
Self-attention constructs query, key, and value matrices through learned
linear projections:

\begin{equation}
\mathbf{Q}=\mathbf{X}\mathbf{W}^{Q}, \qquad
\mathbf{K}=\mathbf{X}\mathbf{W}^{K}, \qquad
\mathbf{V}=\mathbf{X}\mathbf{W}^{V}.
\label{eq:qkv_projections}
\end{equation}
Here, $\mathbf{W}^{Q}$, $\mathbf{W}^{K}$, and $\mathbf{W}^{V}$ are
learned projection matrices, while $\mathbf{Q}$, $\mathbf{K}$, and
$\mathbf{V}$ denote the resulting query, key, and value matrices.

The query associated with one position is compared with the keys associated
with all positions. Scaled dot-product attention is then calculated as

\begin{equation}
\operatorname{Attention}(\mathbf{Q},\mathbf{K},\mathbf{V})
=
\operatorname{softmax}
\left(
\frac{\mathbf{Q}\mathbf{K}^{\mathsf{T}}}{\sqrt{d_k}}
\right)\mathbf{V},
\label{eq:scaled_dot_product_attention}
\end{equation}

where $d_k$ is the dimensionality of the key vectors
\cite{vaswani2017attention}. The scaling factor $\sqrt{d_k}$ limits the
magnitude of the dot products before the softmax operation. Each row of the
resulting attention matrix defines a distribution over the positions from
which a query receives information.

This mechanism provides a short computational path between any pair of
positions. However, self-attention without additional positional information
is insensitive to the order in which the input representations are supplied:
permuting the input positions produces the corresponding permutation of the
outputs. Consequently, the Transformer requires either explicit positional
representations or an architectural mechanism through which position can
affect the attention calculation. This requirement is central to the
length-generalisation problem considered in this dissertation.

\subsection{Multi-Head Attention and Transformer Encoder Blocks}
\label{subsec:multihead_encoder}

Multi-head attention applies the attention calculation using multiple sets of
learned projections. For attention head $h$,

\begin{equation}
\operatorname{head}_h
=
\operatorname{Attention}
\left(
\mathbf{X}\mathbf{W}^{Q}_h,
\mathbf{X}\mathbf{W}^{K}_h,
\mathbf{X}\mathbf{W}^{V}_h
\right).
\label{eq:attention_head}
\end{equation}

The head outputs are concatenated and projected back to the model dimension:

\begin{equation}
\operatorname{MultiHead}(\mathbf{X})
=
\operatorname{Concat}
\left(
\operatorname{head}_1,\ldots,\operatorname{head}_{N{\mathrm{head}}}
\right)\mathbf{W}^{O},
\label{eq:multihead_attention}
\end{equation}

where $N{\mathrm{head}}$ is the number of attention heads and $\mathbf{W}^{O}$ is a learned
output projection. Using separate projections allows different heads to form
different attention distributions, although the existence of multiple heads
does not guarantee that they learn distinct or interpretable functions.

A Transformer encoder block combines multi-head self-attention with a
position-wise feed-forward network. Residual connections preserve the input
to each sublayer, while layer normalisation is applied to stabilise the
representations. In the post-normalisation formulation used by the original
Transformer, the two principal operations can be written as

\begin{align}
\mathbf{H}
&=
\operatorname{LayerNorm}
\left(
\mathbf{X}+\operatorname{MultiHead}(\mathbf{X})
\right), \\
\mathbf{Z}
&=
\operatorname{LayerNorm}
\left(
\mathbf{H}+\operatorname{FFN}(\mathbf{H})
\right),
\label{eq:transformer_encoder_block}
\end{align}

with

\begin{equation}
\operatorname{FFN}(\mathbf{H})
=
\sigma(\mathbf{H}\mathbf{W}_1+\mathbf{b}_1)\mathbf{W}_2+\mathbf{b}_2,
\label{eq:positionwise_ffn}
\end{equation}

Here, $\mathbf{H}$ is the representation produced by the attention
sublayer and $\mathbf{Z}$ is the output of the encoder block.
$\operatorname{LayerNorm}$ denotes layer normalisation, while
$\operatorname{FFN}$ denotes the position-wise feed-forward network.
The matrices $\mathbf{W}_1$ and $\mathbf{W}_2$ and the vectors
$\mathbf{b}_1$ and $\mathbf{b}_2$ are learned parameters, and $\sigma$
is a non-linear activation function. The same feed-forward transformation
is applied independently at every sequence position.

The synthetic models used in this dissertation employ compact stacks of
encoder blocks following this general structure. Their precise dimensions,
activation function, optimisation settings, and output mappings are reported
in Chapter~\ref{chap:Methodology and Experimental Design}. The next section examines the mechanisms
through which positional information is incorporated into these blocks.

\section{Positional Information in Transformers}
\label{sec:positional_background}

Unlike recurrent models, self-attention does not process tokens through an
inherently ordered sequence of state updates. Without an additional source of
positional information, the same token representations presented in a
different order produce correspondingly permuted outputs. A Transformer must
therefore obtain information about order either through explicit position
vectors, modifications to the attention calculation, or structural features
such as an autoregressive mask.

Positional methods can broadly be distinguished by whether they represent
absolute positions or relationships between positions. Absolute methods
associate a representation with each position in the sequence. Relative
methods instead modify attention according to the distance or displacement
between a query and a key \cite{shaw2018relative}. This distinction is
particularly relevant to length generalisation because a model may be
evaluated at absolute positions or relative distances that were not observed
during training.

\subsection{Learned and Sinusoidal Positional Encodings}
\label{subsec:absolute_positional_encodings}

Learned absolute positional embeddings associate each position $p$ with a
trainable vector $\mathbf{p}_p$. The input representation at that position is

\begin{equation}
\mathbf{z}_p = \mathbf{e}(x_p) + \mathbf{p}_p,
\label{eq:learned_position_embedding}
\end{equation}

where $\mathbf{e}(x_p)$ is the embedding of the token at position $p$.
Because the positional vectors are optimised jointly with the model, they can
adapt to the positions represented in the training data. However, this
approach does not impose a functional relationship between neighbouring
positions. Entries allocated for positions beyond the training range receive
no direct training signal, and positions outside the allocated embedding
table cannot be represented without extending the table.

The original Transformer also introduced fixed sinusoidal positional
encodings \cite{vaswani2017attention}. For position $p$ and representation
dimension $i$, these are defined as

\begin{align}
\operatorname{PE}(p,2i)
&=
\sin\left(
\frac{p}{10000^{2i/d_{\mathrm{model}}}}
\right), \\
\operatorname{PE}(p,2i+1)
&=
\cos\left(
\frac{p}{10000^{2i/d_{\mathrm{model}}}}
\right).
\label{eq:sinusoidal_position_encoding}
\end{align}

Here, $\operatorname{PE}$ denotes the positional encoding, $p$ is the sequence position, and $i$ indexes pairs of representation dimensions.

Different dimensions therefore vary at different frequencies. Unlike a
learned lookup table, the sinusoidal representation is generated by a fixed
function and can be evaluated at positions beyond those used during
training. Relative offsets can also be expressed through linear
transformations of these sinusoidal components. These properties make
sinusoidal encoding a plausible candidate for length extrapolation, but they
do not guarantee that the surrounding Transformer learns a computation that
uses the representation systematically at unseen lengths.

\subsection{Rotary Position Embedding}
\label{subsec:rope_background}

Rotary Position Embedding (RoPE) incorporates position into the query and key
vectors rather than adding a positional vector to the token representation
\cite{su2021roformer}. Pairs of query and key dimensions are rotated by an
angle determined by their position. For position $p$, the transformation can
be written abstractly as

\begin{equation}
\widetilde{\mathbf{q}}_p
=
\mathbf{R}_{p}\mathbf{q}_p,
\qquad
\widetilde{\mathbf{k}}_p
=
\mathbf{R}_{p}\mathbf{k}_p,
\label{eq:rope_rotation}
\end{equation}

where $\mathbf{R}_{p}$ is a block-diagonal rotation matrix containing
position-dependent sine and cosine terms. Attention is subsequently computed
using the rotated query and key vectors:

\begin{equation}
\widetilde{\mathbf{q}}_p^{\mathsf{T}}
\widetilde{\mathbf{k}}_r
=
\mathbf{q}_p^{\mathsf{T}}
\mathbf{R}_{p}^{\mathsf{T}}
\mathbf{R}_{r}
\mathbf{k}_r.
\label{eq:rope_relative}
\end{equation}

Because the product
$\mathbf{R}_{p}^{\mathsf{T}}\mathbf{R}_{r}$ depends on the relative
displacement between $p$ and $r$, RoPE introduces relative-position
information into the attention score while retaining an absolute
position-dependent transformation. The rotations can be evaluated at longer
positions, although their mathematical availability does not ensure that a
model trained over a shorter range will interpret the unseen rotation phases
correctly.

\subsection{Attention with Linear Biases}
\label{subsec:alibi_background}

Attention with Linear Biases (ALiBi) represents distance by adding a
head-specific linear penalty directly to each attention score
\cite{press2022train}. A general form of the biased score between positions
$i$ and $j$ in head $h$ is

\begin{equation}
a_{ij}^{(h)}
=
\frac{
\mathbf{q}_{i}^{(h)}
\left(\mathbf{k}_{j}^{(h)}\right)^{\mathsf{T}}
}{
\sqrt{d_k}
}
-
m_h\Delta(i,j),
\label{eq:alibi_attention}
\end{equation}

where $m_h$ is a fixed slope and $\Delta(i,j)$ is the positional distance
used by the attention mechanism. Different slopes cause different heads to
apply different distance penalties. In a bidirectional encoder,
$\Delta(i,j)$ can be defined using absolute distance, whereas a causal model
restricts attention to preceding positions.

ALiBi does not require a positional embedding table and its bias can be
calculated at arbitrary sequence lengths. It was specifically introduced
under a train-short/test-long evaluation setting. Nevertheless, the linear
bias favours particular attention-distance patterns; whether those patterns
support extrapolation depends on the task and on the dependencies required
to produce the correct output.

\subsection{Transformers Without Explicit Positional Encoding}
\label{subsec:nope_background}

A Transformer without explicit positional encoding is commonly described as
NoPE. Under this condition, no positional vector is added to the token
embedding, and neither RoPE nor a positional attention bias is applied.
Research on causal Transformers has shown that models can sometimes recover
useful position-related behaviour without an explicit positional encoding,
partly because the causal attention mask itself introduces an ordering
structure \cite{kazemnejad2023impact}.

This observation does not mean that explicit position information is
unnecessary in every architecture. In an unmasked bidirectional
self-attention encoder, removing all positional mechanisms leaves the model
permutation-equivariant: reordering the input tokens causes the outputs to be
reordered in the same way, without otherwise identifying their original
positions. NoPE is consequently an important control condition. It tests
whether task content and the remaining architecture provide enough
information to learn the required transformation.

The five settings considered in this dissertation therefore introduce
position at different points in the computation. Learned and sinusoidal
methods modify the input representations; RoPE transforms queries and keys;
ALiBi modifies the attention scores; and NoPE supplies no explicit position
signal. Their structural differences motivate the controlled comparison
reported in the later experimental chapters.

\section{Length Generalisation}
\label{sec:length_generalisation_background}

\subsection{Length Generalisation and Out-of-Distribution Evaluation}
\label{subsec:length_ood}

Consider a family of tasks in which $f_n$ denotes the required mapping for
sequences of length $n$. A model is trained using examples whose lengths
belong to a training set $\mathcal{L}_{\mathrm{train}}$, with maximum observed
length

\begin{equation}
L_{\max}
=
\max \mathcal{L}_{\mathrm{train}}.
\label{eq:maximum_training_length}
\end{equation}

An ordinary in-distribution evaluation measures performance on new examples
whose lengths remain within $\mathcal{L}_{\mathrm{train}}$. In contrast, a
length-generalisation evaluation uses test lengths

\begin{equation}
n > L_{\max},
\label{eq:unseen_test_length}
\end{equation}

while preserving the underlying task definition. The inputs are therefore
longer, but the rule required to produce the output remains unchanged.

This setting differs from standard independent and identically distributed
testing. The model must extrapolate along a controlled dimension of the data
rather than generalise only to new samples drawn from the same length
distribution. Strong performance at $L_{\max}$ demonstrates that the model
can fit the observed task instances, but it does not establish that the
learned computation represents the complete family $\{f_n\}$. A solution can
be correct over the observed range while depending on length-specific
features, absolute positions, or statistical regularities that cease to hold
when $n>L_{\max}$.

Length generalisation is related to systematic generalisation: both concern
whether learned components or operations can be applied in combinations
outside the training distribution. Work on compositional sequence learning
has shown that high performance on familiar combinations can coexist with
poor performance under systematic changes to the test distribution
\cite{lake2018generalization}. Length extrapolation provides a particularly
controlled instance of this broader problem because the task semantics can be
held constant while only the size of the instance changes.

A rigorous train-short/test-long experiment must therefore separate two
questions. First, did the model reliably learn the task at the training
length? Second, conditional on that success, did the learned solution transfer
to longer instances? Without the first condition, poor long-length accuracy
cannot be attributed specifically to extrapolation.

\subsection{Length Generalisation in Algorithmic Tasks}
\label{subsec:algorithmic_length_generalisation}

Algorithmic tasks provide a controlled environment for studying this form of
generalisation. Their target outputs are produced by deterministic procedures,
and new examples can be generated at precisely specified lengths. In
principle, an algorithm such as digit-wise addition or sequence reversal can
be applied to inputs longer than those used during training. In practice,
however, a neural model can minimise its training loss using a representation
that is only reliable over the observed range.

This distinction has motivated research into whether successful neural
models learn transferable procedures or length-specific approximations.
Studies of arithmetic with Transformer-based language models report that
increasing model scale alone does not reliably produce length extrapolation,
while intermediate reasoning steps can improve some settings
\cite{anil2022exploring}. More broadly, experiments organised according to
formal-language complexity show that fitting the training distribution does
not imply non-trivial generalisation to longer strings, and that performance
depends on both the task structure and the inductive bias of the architecture
\cite{deletang2023neural}.

Zhou et al.\ examine this question through the algorithms representable by
Transformers and distinguish tasks for which training tends to recover a
simple length-generalising program from tasks for which models favour
non-generalising pattern-matching solutions \cite{zhou2024algorithms}. Their
results are important because they show that architectural expressivity is
not sufficient by itself: the existence of a Transformer implementation of
an algorithm does not guarantee that optimisation will recover that
implementation from finite training examples.

Synthetic algorithmic benchmarks are therefore useful not because they
represent the full complexity of practical applications, but because they
make the required rule and the distribution shift explicit. They permit
controlled comparisons across tasks, encodings, training procedures, and
random seeds. Their simplicity also makes it possible to analyse individual
output positions and task-specific dependencies after aggregate accuracy has
deteriorated.

\subsection{Length Generalisation Versus Long-Context Modelling}
\label{subsec:length_vs_long_context}

Length generalisation should not be treated as synonymous with long-context
modelling. A long-context architecture is designed to process or retain
information across a large input. Transformer-XL, for example, extends the
effective context available to a model through segment-level recurrence and a
relative positional formulation \cite{dai2019transformerxl}. The Long Range
Arena benchmark instead evaluates the efficiency and effectiveness of
architectures on several tasks containing long sequences
\cite{tay2021longrange}.

These approaches address important computational and representational
questions, but operating on a long sequence does not necessarily demonstrate
extrapolation. A model can be trained and tested on similarly long examples,
in which case the evaluation remains within the observed length regime.
Conversely, a compact model can be used in a strict length-generalisation
experiment even when its test sequences are modest by contemporary
long-context standards, provided that they extend substantially beyond the
training range.

The two research areas can nevertheless interact. Positional mechanisms,
attention-distance biases, memory structures, and computational cost all
influence how a model behaves as its input grows. The defining feature of a
length-generalisation experiment is therefore not the absolute magnitude of
the context window, but the controlled separation between the lengths used
for training and those used for evaluation.

\section{Controlled Sequence Tasks}
\label{sec:controlled_tasks_background}

Controlled sequence tasks make it possible to vary sequence length without
changing the rule that determines the correct output. Unlike natural-language
datasets, their target mappings are deterministic and their relevant
dependencies can be stated explicitly. This makes them useful for separating
failure to learn a task from failure to apply a learned computation at longer
lengths.

The tasks used in this dissertation represent three different computational
requirements. Addition requires symbolic computation with possible carry
propagation; delayed copy requires retention and aligned retrieval after an
intervening delay; and reverse requires a position-dependent reordering of the
complete input. The following subsections describe the computational
properties of these tasks. Their precise data-generation and serialisation
procedures are specified in Chapter~\ref{chap:methodology}.

\subsection{Addition}
\label{subsec:addition_background}

Addition is frequently used as a controlled test of whether neural models can
learn arithmetic procedures rather than associations limited to the numerical
range represented during training \cite{anil2022exploring,
lee2024teaching}. Consider two $n$-digit decimal operands with digit
sequences

\begin{equation}
\mathbf{a}=(a_1,\ldots,a_n),
\qquad
\mathbf{b}=(b_1,\ldots,b_n),
\label{eq:addition_operands}
\end{equation}

where $a_1$ and $b_1$ are the most significant digits. The corresponding
integers are

\begin{equation}
A=\sum_{i=1}^{n}a_i10^{n-i},
\qquad
B=\sum_{i=1}^{n}b_i10^{n-i}.
\label{eq:addition_integer_values}
\end{equation}

Digit-wise addition proceeds from the least significant position towards the
most significant position. If $c_{i+1}$ denotes the carry entering position
$i$, then

\begin{align}
s_i &= a_i+b_i+c_{i+1}, \\
y_i &= s_i \bmod 10, \\
c_i &= \left\lfloor\frac{s_i}{10}\right\rfloor,
\label{eq:digitwise_addition}
\end{align}

with $c_{n+1}=0$. Here, $y_i$ is the output digit at position $i$, while
$c_i$ is the carry propagated towards the next more significant position.

The task therefore combines aligned local interactions with a dependency that
can propagate across several digit positions. A model may perform well on
short operands by learning statistical or position-specific regularities
without recovering a procedure that remains reliable for longer operands.
Furthermore, aggregate accuracy alone does not reveal whether failures are
associated with carry propagation, output alignment, or a broader collapse
across all positions. These alternatives motivate the carry-conditioned and
position-wise diagnostics used later in the dissertation.

\subsection{Delayed Copy}
\label{subsec:copy_background}

Copying tasks have been used to evaluate whether sequence models can preserve
information across an intervening delay \cite{arjovsky2016unitary}. Given a
symbol sequence

\begin{equation}
\mathbf{x}=(x_1,\ldots,x_n),
\label{eq:copy_source}
\end{equation}

the delayed-copy task first presents the sequence, followed by a separator and
a region of padding symbols. The required output reproduces the original
symbols after this delay. Abstractly, the input and target can be represented
as

\begin{align}
\mathbf{x}_{\mathrm{input}}
&=
(x_1,\ldots,x_n,\mathtt{SEP},
\underbrace{\mathtt{PAD},\ldots,\mathtt{PAD}}_{n}), \\
\mathbf{x}_{\mathrm{target}}
&=
(\underbrace{\mathtt{PAD},\ldots,\mathtt{PAD}}_{n+1},
x_1,\ldots,x_n).
\label{eq:delayed_copy_mapping}
\end{align}

Here, $\mathtt{SEP}$ marks the boundary between the source and delayed-output
regions, and $\mathtt{PAD}$ is a padding symbol. Each predicted content token
must match a corresponding token from the first part of the input.

Increasing $n$ simultaneously increases the amount of information that must
be retained and the distance between each source symbol and its required
output position. The task does not require arithmetic transformation, but it
does require reliable content preservation and positional alignment. It
therefore provides a controlled contrast with addition: errors can indicate a
failure to retrieve or align stored symbols rather than a failure of symbolic
calculation.

\subsection{Reverse}
\label{subsec:reverse_background}

For the reverse task, an input sequence

\begin{equation}
\mathbf{x}=(x_1,x_2,\ldots,x_n)
\end{equation}

is mapped to

\begin{equation}
\mathbf{y}=(x_n,x_{n-1},\ldots,x_1),
\qquad
y_i=x_{n-i+1}.
\label{eq:reverse_mapping}
\end{equation}

Unlike delayed copy, reverse does not preserve the original positional
alignment. Every output position must identify the input position determined
by the sequence length. The required source position for a fixed output index
therefore changes when $n$ changes. This makes reverse a direct test of
whether a model has learned a general position-dependent transformation
rather than associations tied to absolute training positions.

The source-target dependency distance also varies across the output. Tokens
near the ends of the sequence must be transferred across large distances,
whereas positions near the centre have shorter dependencies. This structure
allows performance to be analysed according to dependency distance rather
than only output position. Prior controlled studies have similarly used
formal and algorithmic sequence transformations to examine whether neural
architectures transfer learned computations to longer inputs
\cite{deletang2023neural,zhou2024algorithms}.

Figure~\ref{fig:controlled_task_examples} summarises representative
input--target mappings for the three controlled tasks.

\begin{figure}[tbp]
    \centering
    \includegraphics[width=\textwidth]
    {controlled_task_examples.pdf}
    \caption{Representative input--target mappings for addition, delayed
    copy, and reverse. Addition combines digit-wise computation with carry
    propagation; delayed copy requires retrieval after an intervening padded
    region; and reverse maps each output position to its length-dependent
    source position. \texttt{PAD} and \texttt{SEP} denote padding and
    separator tokens, respectively. The examples illustrate task
    serialisation rather than the sequence lengths used in the experiments.}
    \label{fig:controlled_task_examples}
\end{figure}

Together, the three tasks separate computational properties that would be
conflated in a single aggregate benchmark. Addition emphasises computation
and carry propagation, delayed copy emphasises retention and aligned
retrieval, and reverse emphasises length-dependent positional transformation.
This separation supports a failure analysis that asks not only whether
accuracy decreases, but whether different task structures produce different
error behaviours.

\section{Evaluating Length-Generalisation Failure}
\label{sec:evaluation_background}

A length-generalisation evaluation must establish more than whether aggregate
accuracy decreases. It must first verify that the model learned the
training-length task, then identify where deterioration begins and how errors
are distributed. Sequence-level accuracy, token-level accuracy, random-seed
variation, position-wise profiles, and attention statistics provide
complementary evidence about these behaviours.

\subsection{Sequence-Level and Token-Level Metrics}
\label{subsec:sequence_token_metrics}

Let $N_L$ denote the number of evaluation examples at test length $L$.
After applying any task-specific padding mask, let $\mathbf{y}^{(r)}$ and
$\widehat{\mathbf{y}}^{(r)}$ denote the evaluated target and predicted token
sequences for example $r$. Exact-match accuracy is

\begin{equation}
A_{\mathrm{EM}}(L)
=
\frac{100}{N_L}
\sum_{r=1}^{N_L}
\mathbb{I}
\left[
\widehat{\mathbf{y}}^{(r)}
=
\mathbf{y}^{(r)}
\right],
\label{eq:exact_match_accuracy}
\end{equation}

where $\mathbb{I}[\cdot]$ is the indicator function. An example contributes
to the exact-match score only when every predicted token is correct. Exact
match therefore measures complete task success, but it does not show how many
tokens were incorrect within a failed sequence.

If target example $r$ contains $T_r$ evaluated tokens, token accuracy is

\begin{equation}
A_{\mathrm{tok}}(L)
=
\frac{
100
}{
\sum_{r=1}^{N_L}T_r
}
\sum_{r=1}^{N_L}
\sum_{j=1}^{T_r}
\mathbb{I}
\left[
\widehat{y}^{(r)}_j
=
y^{(r)}_j
\right].
\label{eq:token_accuracy}
\end{equation}

Token accuracy distinguishes almost-correct predictions from predictions that
are incorrect across nearly every output position. The corresponding error
percentages are

\begin{align}
E_{\mathrm{EM}}(L)
&=
100-A_{\mathrm{EM}}(L), \\
E_{\mathrm{tok}}(L)
&=
100-A_{\mathrm{tok}}(L).
\label{eq:error_percentages}
\end{align}

Reporting errors as percentages makes the magnitude of failure directly
interpretable. For example, zero exact-match accuracy always corresponds to
$100\%$ exact-match error, but the token-error percentage may range from a
small value to a value close to uniform prediction.

Exact match also becomes increasingly strict as the output grows. If token
predictions were independent and each token were correct with probability
$p$, the probability of an entirely correct sequence of length $T$ would be

\begin{equation}
P(\text{exact match})=p^T.
\label{eq:exact_match_length_effect}
\end{equation}

This expression is an illustration rather than an independence assumption
about the models. It shows why sequence-level performance can fall sharply
even when token accuracy decreases more gradually.

For a target vocabulary
$\mathcal{V}_{\mathrm{target}}$, uniform prediction gives expected token
accuracy $100/|\mathcal{V}_{\mathrm{target}}|$. The corresponding token-error
reference is

\begin{equation}
E_{\mathrm{chance}}
=
100
\left(
1-\frac{1}{|\mathcal{V}_{\mathrm{target}}|}
\right).
\label{eq:uniform_error_reference}
\end{equation}

This value provides a descriptive reference for near-uniform prediction; it
is not, by itself, a statistical significance test.

\subsection{Random-Seed Stability}
\label{subsec:seed_stability_background}

Neural-network training is stochastic because parameter initialisation, data
ordering, and optimisation can vary between runs. Reporting only the
best-performing seed can therefore conceal unstable training behaviour.
Empirical studies have demonstrated that random seeds can materially change
the conclusions drawn from sequence-model experiments, motivating the
reporting of score distributions or multiple-run summaries
\cite{reimers2017reporting}.

For a metric $A_s(L)$ obtained from seed $s$, the mean across $S$ seeds is

\begin{equation}
\overline{A}(L)
=
\frac{1}{S}
\sum_{s=1}^{S}A_s(L),
\label{eq:seed_mean}
\end{equation}

and the sample standard deviation is

\begin{equation}
\operatorname{SD}(L)
=
\sqrt{
\frac{1}{S-1}
\sum_{s=1}^{S}
\left(
A_s(L)-\overline{A}(L)
\right)^2
}.
\label{eq:seed_standard_deviation}
\end{equation}

In this dissertation, error bars representing
$\overline{A}(L)\pm 1\operatorname{SD}(L)$ describe variation across random
seeds. They do not represent standard errors or confidence intervals. Because
the main benchmark uses three seeds, individual seed outcomes are also
retained when assessing whether a configuration learned the training task
reliably.

This distinction supports the baseline-eligibility procedure introduced in
Chapter~\ref{chap:methodology}. A configuration that succeeds for only one
initialisation is evidence of seed instability, whereas a configuration that
succeeds consistently at the training length but fails at longer lengths
provides a valid basis for studying extrapolation.

\subsection{Position-Wise and Task-Specific Diagnostics}
\label{subsec:positionwise_background}

Aggregate token error can conceal substantial variation across an output.
Let $N_j$ be the number of evaluated examples containing output position $j$.
The position-wise token-error percentage is

\begin{equation}
E_j
=
100
\left[
1-
\frac{1}{N_j}
\sum_{r=1}^{N_j}
\mathbb{I}
\left[
\widehat{y}^{(r)}_j=y^{(r)}_j
\right]
\right].
\label{eq:positionwise_error}
\end{equation}

Plotting $E_j$ across positions reveals whether errors are approximately
uniform, concentrated near sequence boundaries, or progressively increasing
through the output. Positions can also be grouped into start, middle, and end
regions to provide compact percentage summaries without discarding the full
position-wise profile.

Task-specific groupings test more focused hypotheses. For addition, output
digits can be divided according to whether they receive an incoming carry.
For reverse, errors can be grouped by the distance between each output
position and its required source position. Delayed copy can be examined
according to the aligned source and target locations. These analyses identify
associations between error and task structure, although they do not alone
establish the internal cause of a prediction.

\subsection{Attention-Based Diagnostics}
\label{subsec:attention_diagnostics_background}

Attention weights provide an additional description of how the model
distributes attention mass. Let $\alpha_{ij}^{(h)}$ denote the attention
weight from query position $i$ to key position $j$ in head $h$, with

\begin{equation}
\sum_{j=1}^{n}\alpha_{ij}^{(h)}=1.
\end{equation}

The normalised attention entropy is

\begin{equation}
H_i^{(h)}
=
-
\frac{
\sum_{j=1}^{n}
\alpha_{ij}^{(h)}
\log\alpha_{ij}^{(h)}
}{
\log n
}.
\label{eq:normalised_attention_entropy}
\end{equation}

Values closer to zero indicate concentrated attention, while values closer to
one indicate a more diffuse distribution. Dividing by $\log n$ makes entropy
more comparable across sequence lengths.

The average attention distance is

\begin{equation}
D_i^{(h)}
=
\sum_{j=1}^{n}
\alpha_{ij}^{(h)}
|i-j|,
\label{eq:average_attention_distance}
\end{equation}

and the local-attention ratio within a window of radius $w$ is

\begin{equation}
R_i^{(h)}(w)
=
\sum_{\substack{j=1\\|i-j|\leq w}}^{n}
\alpha_{ij}^{(h)}.
\label{eq:local_attention_ratio}
\end{equation}

These quantities describe whether attention becomes more diffuse, operates
over greater distances, or remains concentrated near each query position as
sequence length increases. They must nevertheless be interpreted alongside
accuracy. Average distance can grow mechanically when longer positions become
available, and a change in attention statistics does not necessarily identify
the cause of a model's error.

Attention weights are therefore treated in this dissertation as behavioural
diagnostics rather than complete causal explanations. Prior research has
shown that substantially different attention distributions can sometimes
produce similar predictions, cautioning against interpreting attention
weights as faithful explanations of model decisions
\cite{jain2019attention}. Their role here is to identify changes associated
with extrapolation failure, not to claim that an observed attention pattern
alone caused that failure.

\section{Training Strategies for Length Generalisation}
\label{sec:mitigation_background}

Training interventions for length generalisation address different possible
sources of failure. Exposing a model to several training lengths targets
dependence on a single length distribution. Curriculum learning targets the
optimisation path by increasing task difficulty gradually. Randomised
positional exposure instead targets the possibility that the positional
representations encountered at longer lengths are out of distribution.

These strategies should not be interpreted as interchangeable. An
intervention may improve optimisation without changing extrapolation, or it
may expose additional positions while weakening performance on the original
training task. Its intended target and its training-length performance must
therefore be established before unseen-length results are interpreted.

\subsection{Mixed-Length Training}
\label{subsec:mixed_length_background}

A fixed-length training condition draws every training example from a single
sequence length $L_{\mathrm{train}}$. Mixed-length training instead draws
examples from a set

\begin{equation}
\mathcal{L}_{\mathrm{mix}}
=
\{L_1,L_2,\ldots,L_K\},
\label{eq:mixed_length_set}
\end{equation}

where the lengths are no greater than the maximum training length. If
$\pi_k$ is the proportion of training examples drawn from length $L_k$, the
mixed-length objective can be written as

\begin{equation}
\mathcal{J}_{\mathrm{mix}}(\boldsymbol{\theta})
=
\sum_{k=1}^{K}
\pi_k
\,
\mathbb{E}_{
(\mathbf{x},\mathbf{y})\sim\mathcal{D}_{L_k}
}
\left[
\ell
\left(
f_{\boldsymbol{\theta}}(\mathbf{x}),
\mathbf{y}
\right)
\right],
\label{eq:mixed_length_objective}
\end{equation}

where $\boldsymbol{\theta}$ denotes the model parameters,
$\mathcal{D}_{L_k}$ is the data distribution at length $L_k$, and
$\ell$ is the training loss.

The hypothesis is that exposure to multiple lengths discourages a model from
associating the task with one fixed sequence shape or one set of boundary
positions. It may instead encourage features that are shared across the
observed length range. This motivation is consistent with evidence that
Transformer length generalisation depends on both the task and the training
distribution \cite{anil2022exploring,zhou2024algorithms}.

Mixed-length training introduces a trade-off. For a fixed total data budget,
including more lengths reduces the number of examples available at each
individual length. Padding and batching can also change the proportion of
content and special tokens seen during optimisation. Consequently, a
mixed-length model may underfit even the original training task. Any apparent
extrapolation benefit must therefore be considered together with
training-length accuracy and a matched fixed-length baseline.

\subsection{Curriculum Learning}
\label{subsec:curriculum_background}

Curriculum learning presents training examples in a progression from simpler
to more difficult conditions \cite{bengio2009curriculum}. In a
length-based curriculum, sequence length acts as the difficulty variable. A
model can be trained through an ordered series

\begin{equation}
L_1 < L_2 < \cdots < L_K,
\label{eq:length_curriculum}
\end{equation}

where the parameters learned at stage $k-1$ initialise training at stage
$k$. Abstractly,

\begin{equation}
\boldsymbol{\theta}_k
=
\operatorname{Train}
\left(
\boldsymbol{\theta}_{k-1},
\mathcal{D}_{L_k}
\right).
\label{eq:curriculum_update}
\end{equation}

The intended target is optimisation. Shorter examples may allow the model to
learn a partial or simpler form of the required computation before longer
dependencies are introduced. The later stages then refine the existing
parameters rather than learning the longest training task entirely from a
random initialisation.

A successful curriculum does not necessarily produce extrapolation. The
model may retain useful features from earlier stages, but it may also adapt to
the final curriculum length and remain dependent on that observed range.
Curriculum learning must therefore be compared with a model trained directly
at the final training length. The comparison tests whether gradual length
increases alter the resulting behaviour, rather than assuming that the
curriculum is intrinsically beneficial.

\subsection{Randomised Positional Exposure and Padded Training}
\label{subsec:randomised_position_background}

A separate hypothesis is that failure occurs because positional
representations at unseen lengths differ from those encountered during
training. Ruoss et al.\ propose randomised positional encodings that simulate
a larger positional range during training
\cite{ruoss2023randomized}. For a sequence of length $n$, their approach
selects an ordered subset

\begin{equation}
p_1 < p_2 < \cdots < p_n
\end{equation}

This formulation describes the randomised positional-index method proposed
by Ruoss et al.; it is not the padded mixed-length implementation evaluated
in this dissertation.

from a larger range of possible indices and applies the positional
representation associated with each selected $p_i$. This allows short
training sequences to encounter positional indices that would ordinarily
appear only in longer inputs. Their method directly targets positional
out-of-distribution exposure.

Padding-based strategies are related but not equivalent. Presenting examples
of several logical lengths within a common padded tensor can vary the
locations of separators, output regions, and sequence boundaries. However,
right-padding alone does not randomly change the starting position of the
content, nor does it reproduce randomised positional-index sampling.
Padding-based training should consequently be described according to the
positions it actually varies.

The padded condition evaluated in this dissertation is motivated by the
broader positional-exposure hypothesis, but is treated as a distinct
intervention rather than a reproduction of randomised positional encoding.
Its precise construction, including the use of mixed logical lengths,
right-padding, and shuffled training examples, is described in
Chapter~\ref{chap:Methodology and Experimental Design

\subsection{Intervention-Specific Evaluation}
\label{subsec:intervention_evaluation}

The effect of an intervention at length $L$ can be expressed relative to its
matched original baseline as

\begin{equation}
\Delta A(L)
=
A_{\mathrm{intervention}}(L)
-
A_{\mathrm{baseline}}(L).
\label{eq:intervention_difference}
\end{equation}

A positive value at unseen lengths is meaningful only if the intervention
also preserves adequate performance at the training length. If the
intervention fails to learn the training task, the result indicates
training-length underfitting rather than evidence about improved or worsened
extrapolation.

The three strategy families therefore test different hypotheses:
mixed-length training tests whether broader length exposure reduces
dependence on a single training length; curriculum learning tests whether a
gradual optimisation path produces a more transferable solution; and the
padded mixed-length condition tests whether varying logical boundaries and
output positions within a common tensor changes extrapolation behaviour.
This intervention-specific framing prevents all unsuccessful variants from
being assigned the same interpretation.

\section{Long Sequences in Time-Series Classification}
\label{sec:timeseries_background}

Multivariate time-series classification assigns a single class label to a
sequence containing several measurements at each time step. This differs
fundamentally from the controlled sequence-to-sequence tasks considered in
this dissertation. In Addition, Delayed Copy, and Reverse, increasing the
input length also increases the number of output tokens that must be predicted
correctly. In time-series classification, by contrast, the output remains a
single class label regardless of the input length. Consequently, length-related
behaviour is evaluated using classification accuracy rather than token or
exact-match accuracy.

Transformers are relevant to long time series because self-attention can model
interactions between observations separated by many time steps. However,
long sequences also increase the computational cost of attention and may
contain dependencies at multiple temporal scales. Real-world time series can
introduce additional complications, including irregular sampling, noise, and
variation in the quantity of information available across a sequence.
Specialised approaches such as Rough Transformers therefore use
continuous-time representations and signature-based attention to capture
local and global temporal structure while reducing computational
requirements~\cite{morenopino2024rough}. Rough Transformers are discussed
here as motivation for studying length in real-world time-series models; they
are not implemented in the experiments in this dissertation.

The exploratory real-world experiment uses SelfRegulationSCP2 from the UEA
Multivariate Time Series Classification Archive~\cite{bagnall2018uea}. The
archive was introduced to support reproducible comparison of multivariate
time-series classification methods and provides predefined training and test
splits across datasets from different application domains. SelfRegulationSCP2
provides a useful contrast to the deterministic synthetic tasks because each
example contains multiple real-valued channels and is associated with a
sequence-level class label.

In this dissertation, a compact Transformer is trained using a fixed prefix
of each training series and subsequently evaluated using progressively longer
prefixes. The same five positional-encoding conditions used in the controlled
experiments are considered. This design examines whether the classifier's
performance is sensitive to input lengths beyond the prefix used during
training. Nevertheless, it should not be interpreted as an exact equivalent
of algorithmic length extrapolation. A longer time-series prefix can introduce
additional useful evidence, irrelevant observations, or noise, whereas a
longer synthetic example applies the same deterministic rule to more symbols.

The real-world experiment is therefore treated as a limited exploratory
extension rather than direct validation of the synthetic-task failure
taxonomy. Its purpose is to examine whether length sensitivity also appears
in a noisy sequence-level classification setting and whether its empirical
form differs from the immediate exact-match collapse observed on the
controlled tasks.

\section{Research Gap and Chapter Summary}
\label{sec:background_gap}

\subsection{Research Gap}

Previous work has established that positional representation and training
length can substantially affect Transformer extrapolation. ALiBi was designed
to improve evaluation on sequences longer than those used during
training~\cite{press2022train}, while comparative studies have shown that the
choice of positional encoding can alter length-generalisation
behaviour~\cite{kazemnejad2023impact}. Other work has investigated
algorithmic tasks and randomised positional strategies as ways of improving
generalisation beyond the observed training lengths
~\cite{zhou2024algorithms,ruoss2023randomized}.

However, reporting only accuracy at longer lengths does not fully explain a
model's behaviour. A low out-of-distribution score may result from several
different situations: the model may not have reliably learned the
training-length task, performance may be unstable across random seeds, failure
may occur immediately after the training boundary, or errors may be
concentrated in particular output positions or task-specific conditions.
These cases have different interpretations and should not be grouped together
as a single form of length-generalisation failure.

There is therefore a need for an evaluation framework that separates baseline
validity from extrapolation behaviour. In particular, the training-length
baseline should be verified across multiple seeds before conclusions are
drawn from longer test sequences. Once a valid baseline has been established,
the analysis should examine not only whether accuracy decreases, but also
where errors occur, how close performance is to a task-appropriate chance
reference, and whether the same behaviour is reproduced across positional
encoding choices.

Mitigation experiments require similar care. An intervention can only be
interpreted as improving length generalisation if it first preserves adequate
performance at the training length. Otherwise, low out-of-distribution
performance may simply reflect training underfit. Furthermore, mixed-length
training, curriculum learning, and randomised padding target different
possible causes of failure. Their outcomes should consequently be compared
with matched original baselines and interpreted in relation to the specific
hypothesis each intervention was intended to test.

This dissertation addresses these gaps through a controlled comparison of
five positional-encoding settings across Addition, Delayed Copy, and Reverse.
It combines multi-seed baseline eligibility with exact-match and token-level
metrics, percentage-based error reporting, position-wise diagnostics,
task-specific analyses, and matched baseline-versus-mitigation comparisons.
A limited multivariate time-series classification experiment is included as
an exploratory contrast, while remaining separate from the main synthetic
failure taxonomy.

\subsection{Chapter Summary}

This chapter introduced the Transformer components required for the study,
including self-attention and positional representation. It explained why
performance at the training length is insufficient evidence of length
generalisation and described the distinct computational demands of Addition,
Delayed Copy, and Reverse. It also established the roles of exact-match
accuracy, token error, multi-seed evaluation, position-wise analysis, and
task-specific diagnostics.

The chapter further distinguished the hypotheses underlying mixed-length
training, curriculum learning, and randomised padding, and explained why a
mitigation must retain training-length performance before its extrapolation
behaviour can be interpreted. Finally, it situated the exploratory
SelfRegulationSCP2 experiment within multivariate time-series classification
and clarified why sequence-level classification is not directly equivalent
to deterministic sequence-to-sequence extrapolation.

The next chapter translates these principles into the experimental
methodology, specifying the datasets, model architecture, positional-encoding
conditions, training procedure, baseline-eligibility criterion, evaluation
metrics, diagnostic analyses, and mitigation comparisons.
% -----------------------------------------------------------------------------

\chapter{Methodology and Experimental Design}
\label{chap:Methodology and Experimental Design}

\noindent
This chapter specifies the experimental methodology used to investigate
length generalisation in Transformer-based sequence models. The experiments
follow a controlled train-short/test-long protocol in which a model is trained
using sequences of a fixed length and subsequently evaluated on progressively
longer unseen sequences. This design makes it possible to distinguish
successful learning at the training length from extrapolation beyond the
observed length distribution.

The methodology was organised around three requirements. First, a
configuration had to learn the training-length task reliably before its
longer-length behaviour could be interpreted. Second, conclusions had to be
supported across multiple random seeds rather than inferred from a single
run. Third, evaluation had to describe not only whether performance decreased,
but also where and how the resulting errors were distributed.

The chapter first presents the overall experimental pipeline and then
describes the controlled tasks, dataset generation, Transformer architecture,
positional-encoding conditions, training procedure, baseline-eligibility
criterion, evaluation metrics, diagnostic analyses, and mitigation
experiments. A separate exploratory experiment on multivariate time-series
classification is presented at the end of the chapter. Throughout the
controlled experiments, the architecture and training procedure were kept as
consistent as possible so that sequence length, task, positional encoding,
and training intervention remained the principal experimental variables.


\section{Experimental Overview}
\label{sec:experimental_overview}

The study used a staged experimental pipeline. Controlled algorithmic tasks
were first generated at predefined training and evaluation lengths. A compact
Transformer was then trained for each combination of task, positional
encoding, and random seed. Five positional-encoding conditions were examined:
learned positional embeddings, sinusoidal encoding, no explicit positional
encoding, ALiBi, and RoPE. Each baseline condition was repeated using three
random seeds, giving 45 baseline training runs across the three tasks.

The train-short/test-long protocol follows the general form used in previous
investigations of positional extrapolation and algorithmic length
generalisation~\cite{press2022train,kazemnejad2023impact,zhou2024algorithms}.
However, the present methodology introduces an explicit baseline-eligibility
stage before extrapolation results are interpreted. Training-length
performance was assessed separately for every task--encoding--seed
combination. Only task--encoding configurations that learned the training
task sufficiently reliably across seeds were used as the principal baselines
for the subsequent failure analysis.

Eligible baselines were evaluated over a sequence of increasing test lengths.
The resulting curves were used to determine whether degradation began at the
first unseen length or emerged more gradually. Exact-match accuracy was
supplemented by token-error percentages, position-wise error profiles, and
task-specific diagnostics. These analyses examined whether errors were
distributed uniformly across an output or associated with particular
positions, carry conditions, or dependency distances. Attention analyses were
used as supporting diagnostics rather than as direct evidence that a model
had learned a particular algorithm.

Mitigation experiments were conducted only as matched comparisons with an
original baseline using the same task and positional-encoding condition.
Mixed-length training, curriculum learning, and randomised padding were
treated as separate interventions because they test different hypotheses
about the source of length-related failure. An intervention was not interpreted
as improving extrapolation unless it first retained adequate performance at
the training length.

Finally, a limited experiment using the SelfRegulationSCP2 multivariate
time-series dataset examined input-length sensitivity in a sequence-level
classification setting. This experiment used the same positional-encoding
conditions and multiple random seeds, but it was analysed separately because
time-series classification is not directly equivalent to deterministic
sequence-to-sequence extrapolation.

\begin{table}[htbp]
\centering
\label{tab:experimental_pipeline}
\renewcommand{\arraystretch}{1.2}
\begin{tabular}{
    p{0.08\textwidth}
    p{0.29\textwidth}
    p{0.52\textwidth}
}
\hline
\textbf{Stage} & \textbf{Procedure} & \textbf{Purpose} \\
\hline
1 &
Controlled data generation &
Create deterministic training and test sets at predefined sequence lengths. \\

2 &
Multi-seed baseline training &
Train each task and positional-encoding condition using seeds 42, 123, and
2024. \\

3 &
Baseline eligibility &
Identify configurations that reliably learn the training-length task before
interpreting extrapolation. \\

4 &
Extrapolation diagnostics &
Measure exact-match and token errors across lengths and determine where
failure occurs. \\

5 &
Matched mitigation comparison &
Compare each intervention with its corresponding original baseline and check
whether training-length performance is retained. \\

6 &
Exploratory real-world extension &
Examine input-length sensitivity on SelfRegulationSCP2 using sequence-level
classification accuracy. \\
\hline
\end{tabular}
\caption{Stages of the experimental methodology.}
\end{table}

\section{Synthetic Task Design}
\label{sec:method-synthetic-tasks}

The conceptual motivation for Addition, Delayed Copy, and Reverse was
discussed in Chapter~\ref{chap:background}, and their transformations are
illustrated in Figure~\ref{fig:controlled_task_examples}. This section
therefore focuses on their exact experimental implementation.

For each task, the nominal length \(L\) has a task-specific meaning. In
Addition, \(L\) is the number of digits in each operand; in Delayed Copy, it
is the number of content tokens to reproduce; and in Reverse, it is the
complete sequence length. The resulting sequence construction is summarised
in Table~\ref{tab:synthetic-task-implementation}.

\begin{table}[htbp]
\centering

\label{tab:synthetic-task-implementation}
\small
\renewcommand{\arraystretch}{1.2}
\begin{tabular}{
    p{0.18\textwidth}
    p{0.25\textwidth}
    p{0.26\textwidth}
    p{0.18\textwidth}
}
\hline
\textbf{Task} &
\textbf{Meaning of \(L\)} &
\textbf{Encoded sequence length} &
\textbf{Vocabulary} \\
\hline

Addition &
Digits per operand &
\(2L+1\) input and target tokens &
10 digits, plus, and padding \\

Delayed Copy &
Content tokens to reproduce &
\(2L+1\) input and target tokens &
100 content symbols, separator, and padding \\

Reverse &
Number of input tokens &
\(L\) input and target tokens &
100 content symbols \\

\hline
\end{tabular}
\caption{Implementation of the controlled synthetic tasks.}
\end{table}

For Addition, leading zeros were excluded so that both operands always had
the intended length. The sum was right-aligned within the target and preceded
by padding tokens. For Delayed Copy, the input contained the content sequence,
a separator, and \(L\) padding tokens; the target contained \(L+1\) padding
tokens followed by the original content sequence. Reverse directly mapped
each generated sequence to the same tokens in the opposite order.

Associative Recall was also implemented during development but was excluded
from the main comparison because it did not provide sufficiently reliable
training-length performance. This decision followed the baseline-eligibility
principle: extrapolation behaviour was not interpreted when the original task
had not first been learned reliably.

\section{Train-Short/Test-Long Protocol}
\label{sec:train-short-test-long}

Each model was trained using examples at one nominal length
\(L_{\mathrm{train}}\) and evaluated at both the training length and a sequence
of progressively longer lengths \cite{press2022train,kazemnejad2023impact,zhou2024algorithms}. Performance at
\(L_{\mathrm{test}}=L_{\mathrm{train}}\) was treated as in-distribution
performance, while evaluations satisfying
\(L_{\mathrm{test}}>L_{\mathrm{train}}\) were treated as
out-of-distribution length extrapolation.

Addition models were trained using operands containing 16 digits and evaluated
using operands containing up to 1024 digits. Delayed Copy and Reverse models
were trained at a nominal length of 128 and evaluated at lengths up to 1024.
The complete schedule is shown in
Table~\ref{tab:train-test-lengths}. These values refer to nominal task length:
as specified in Section~\ref{sec:method-synthetic-tasks}, the encoded sequence
length is \(2L+1\) for Addition and Delayed Copy and \(L\) for Reverse.

\begin{table}[htbp]
\centering

\label{tab:train-test-lengths}
\small
\renewcommand{\arraystretch}{1.2}
\begin{tabular}{
    p{0.23\textwidth}
    p{0.15\textwidth}
    p{0.17\textwidth}
    p{0.34\textwidth}
}
\hline
\textbf{Task} &
\(\boldsymbol{L_{\mathrm{train}}}\) &
\textbf{First unseen length} &
\(\boldsymbol{L_{\mathrm{test}}}\) \\
\hline

Addition &
16 &
32 &
16, 32, 64, 128, 256, 512, 1024 \\

Delayed Copy &
128 &
256 &
128, 256, 512, 1024 \\

Reverse &
128 &
256 &
128, 256, 512, 1024 \\

SelfRegulationSCP2 &
256 &
512 &
256, 512, 1024, 1152 \\

\hline
\end{tabular}
\caption{Training and evaluation-length schedule. Synthetic-task values denote
nominal task lengths, while SelfRegulationSCP2 values denote time-series
prefix lengths.}
\end{table}

For the synthetic tasks, a separate test set was generated at every evaluation
length using the same deterministic rule as the training data. The underlying
task therefore remained unchanged; only the length required to execute the
task increased.

Training-length performance was evaluated before extrapolation behaviour. A
configuration that did not learn the task reliably at
\(L_{\mathrm{train}}\) was classified as training underfit and was not used as
evidence of length-generalisation failure. For configurations that satisfied
the baseline-eligibility criterion, performance at the first unseen and
subsequent lengths was used to characterise whether degradation was immediate
or gradual. The numerical eligibility and collapse thresholds are defined in
Section~\ref{sec:evaluation-metrics}.

\section{Experiment Matrix}
\label{sec:experiment-matrix}

The synthetic baseline experiment crossed three tasks, five positional-encoding
conditions, and three random seeds. This produced

\[
3\ \text{tasks}
\times 5\ \text{positional encodings}
\times 3\ \text{seeds}
=45
\]

independently trained synthetic baseline models. Each trained checkpoint was
evaluated at every test length configured for its task. Associative Recall was
not included in this matrix because it did not satisfy the training-fit
requirement described in Section~\ref{sec:method-synthetic-tasks}.

The complete synthetic matrix was first used to assess training-length
performance and seed stability. Baseline eligibility was then determined at
the task--encoding level using results from all three seeds. Consequently,
all 45 baseline runs contributed to the eligibility analysis, while the
principal extrapolation and failure-behaviour analyses focused on
task--encoding configurations that satisfied the eligibility criterion.

The real-world extension crossed one dataset with the same five positional
encodings and three seeds, producing a further 15 trained models. It was
analysed separately from the synthetic benchmark because it predicts one
sequence-level class rather than a target token at every output position.

Table~\ref{tab:experiment-matrix} summarises the baseline experiment matrix.

\begin{table}[htbp]
\centering
\label{tab:experiment-matrix}
\renewcommand{\arraystretch}{1.2}
\begin{tabular}{
    p{0.31\textwidth}
    p{0.48\textwidth}
    p{0.11\textwidth}
}
\hline
\textbf{Experiment component} &
\textbf{Values} &
\textbf{Models} \\
\hline

Synthetic tasks &
Addition, Delayed Copy, Reverse &
-- \\

Synthetic positional encodings &
Learned, sinusoidal, NoPE, ALiBi, RoPE &
-- \\

Synthetic seeds &
42, 123, 2024 &
-- \\

Synthetic baseline matrix &
\(3\) tasks \(\times\) \(5\) encodings \(\times\) \(3\) seeds &
45 \\

Real-world dataset &
SelfRegulationSCP2 &
-- \\

Real-world baseline matrix &
\(1\) dataset \(\times\) \(5\) encodings \(\times\) \(3\) seeds &
15 \\

Total baseline models &
Synthetic and real-world experiments &
60 \\

\hline
\end{tabular}
\caption{Baseline experimental matrix. The counts refer to independently
trained model checkpoints and exclude the subsequent mitigation experiments.}
\end{table}

The task dimension was used to compare behaviour across distinct deterministic
transformations. The positional-encoding dimension examined whether
training fit and extrapolation varied with the representation of sequence
order. Finally, the seed dimension measured sensitivity to model
initialisation, training stochasticity, and synthetic data generation.
Mitigation experiments were conducted as additional matched interventions and
are specified separately in Section~\ref{sec:mitigation-methodology}.

\section{Model Architecture}
\label{sec:model-architecture}

The synthetic experiments used a compact encoder-only Transformer implemented
from scratch in PyTorch, following the principal components of the original
Transformer architecture~\cite{vaswani2017attention}. The model was not
pretrained. It consisted of a task-specific token-embedding layer, one of the
five positional treatments, two Transformer blocks, and a linear output layer
that produced vocabulary logits independently at every sequence position.

Each Transformer block applied multi-head self-attention followed by a
position-wise feedforward network. Residual connections and dropout were
applied around both sublayers, with layer normalisation performed after each
residual addition. The feedforward network used a ReLU activation between its
two linear transformations.

The model used full bidirectional self-attention and did not apply a causal
mask. This was appropriate because every controlled task was formulated as a
non-autoregressive sequence-to-sequence transformation: the complete input
sequence was supplied simultaneously, and the model predicted a target token
at every position. The architecture was held constant across
positional-encoding conditions, except for the positional component being
evaluated.

\begin{table}[htbp]
\centering
\small
\renewcommand{\arraystretch}{1.2}
\begin{tabular}{
    p{0.46\textwidth}
    p{0.38\textwidth}
}
\hline
\textbf{Architecture component} &
\textbf{Configuration} \\
\hline

Token-embedding dimension &
128 \\

Transformer blocks &
2 \\

Attention heads &
4 \\

Dimension per attention head &
32 \\

Feedforward hidden dimension &
256 \\

Feedforward activation &
ReLU \\

Dropout probability &
0.1 \\

Normalisation &
Post-residual layer normalisation \\

Attention visibility &
Full bidirectional attention; no causal mask \\

Output layer &
Linear projection to task vocabulary at every position \\

\hline
\end{tabular}
\caption{Transformer architecture used for the controlled synthetic
experiments. Optimisation settings are reported separately in
Section~\ref{sec:training-procedure}.}
\label{tab:model-config}
\end{table}

The maximum supported sequence length was chosen to cover the longest encoded
input in each task. It was set to 2049 positions for Addition and Delayed Copy,
corresponding to \(2(1024)+1\), and to 1024 positions for Reverse. These
limits determined the size of the learned positional-embedding table and the
precomputed sinusoidal encoding buffer but did not change the number of
Transformer blocks or attention heads.

The architecture was intentionally compact because the objective was
controlled comparison rather than state-of-the-art task performance. Keeping
the architecture fixed made it possible to repeat the full experiment matrix
across tasks, positional encodings, lengths, and seeds while attributing
differences more directly to the experimental conditions.

\section{Positional Encoding Variants}
\label{sec:positional-encoding-variants}

Five positional conditions were evaluated: learned absolute embeddings,
sinusoidal encodings, no explicit positional encoding, ALiBi, and RoPE. Their
theoretical motivations were discussed in Chapter~\ref{chap:background}; this
section specifies how they were implemented in the experimental model.

Only the positional treatment was changed between conditions. The token
embedding, number of Transformer blocks, attention heads, feedforward
dimension, output layer, and optimisation procedure otherwise remained fixed.
This controlled design reduced the possibility that an architectural change
unrelated to position caused the observed performance differences.

\begin{table}[htbp]
\centering
\small
\renewcommand{\arraystretch}{1.2}
\begin{tabular}{
    p{0.17\textwidth}
    p{0.25\textwidth}
    p{0.47\textwidth}
}
\hline
\textbf{Encoding} &
\textbf{Application point} &
\textbf{Experimental implementation} \\
\hline

Learned &
Input representations &
A trainable embedding vector was added for each absolute position. \\

Sinusoidal &
Input representations &
Fixed sine and cosine vectors were added using the original Transformer
frequency schedule. \\

NoPE &
No explicit intervention &
Token embeddings were passed to the first Transformer block without an added
positional signal. \\

ALiBi &
Attention logits &
A head-specific negative bias proportional to the absolute distance between
query and key positions was added before the softmax. \\

RoPE &
Queries and keys &
Pairs of query and key dimensions were rotated using position-dependent sine
and cosine frequencies before their dot product was calculated. \\

\hline
\end{tabular}
\caption{Implementation of the five positional-encoding conditions.}
\label{tab:positional-encodings}
\end{table}

For the learned condition, a trainable positional-embedding table was
allocated up to the maximum supported sequence length. During baseline
training, however, only entries corresponding to positions present at the
training length received gradient updates. Embeddings assigned exclusively to
longer unseen positions therefore remained at their initial values. This is
an inherent property of the evaluated learned absolute-position baseline and
is relevant when interpreting its extrapolation behaviour.

The sinusoidal condition used the fixed sine and cosine construction from the
original Transformer~\cite{vaswani2017attention}. Because these vectors were
defined in advance for every supported position, evaluation at a longer length
did not require learning additional position parameters.

The NoPE condition applied no explicit positional transformation. It was
included because previous work has shown that Transformers without an
explicit positional encoding can still exhibit task-dependent positional and
length-generalisation behaviour~\cite{kazemnejad2023impact}. In the present
bidirectional implementation, NoPE otherwise used the same attention and
feedforward operations as the other conditions.

ALiBi added a head-specific linear penalty to the attention logits
~\cite{press2022train}. Because the model used bidirectional rather than
causal attention, the implemented bias depended on the absolute separation
between positions:

\[
B^{(h)}_{ij}=-m_h |i-j|,
\]

where \(m_h\) is the fixed slope for attention head \(h\). This symmetric
form favoured nearby positions in both directions. It should therefore be
understood as a bidirectional adaptation of the ALiBi distance bias rather
than the original causal-language-model setting.

RoPE applied rotary position information directly to the query and key
representations before attention scores were computed
~\cite{su2024roformer}. The implementation used an even head dimension of 32
and the standard frequency base of \(10\,000\). Value representations were
not rotated.

All five conditions were subjected to the same multi-seed baseline-eligibility
assessment. Consequently, a positional condition that failed to learn a task
at the training length was classified as training underfit rather than as
evidence of extrapolation failure. The eligibility procedure is defined in
Section~\ref{sec:baseline-eligibility}.

\section{Training and Evaluation Procedure}
\label{sec:training-procedure}

All models were implemented and trained using PyTorch
~\cite{paszke2019pytorch}. For each synthetic configuration, the generated
data contained 10,000 training examples, 1,000 validation examples, and 1,000
test examples at every evaluation length. The validation split was generated
as part of the data pipeline but was not used for early stopping or checkpoint
selection in the final experiments.

Each baseline model was trained for 10 epochs. The checkpoint saved after the
final epoch was used for evaluation; no best-validation checkpoint was
selected. Optimisation used AdamW~\cite{loshchilov2019decoupled} with a
learning rate of \(10^{-3}\). No learning-rate scheduler, early stopping, or
gradient clipping was applied. Because the optimiser was instantiated using
the PyTorch defaults apart from the learning rate, its weight-decay coefficient
was \(0.01\).

The loss was token-level cross-entropy. For Addition and Delayed Copy, target
positions containing the padding token were excluded from the loss. Reverse
does not use a padding token, so every target position contributed to its
training loss.

\begin{table}[htbp]
\centering
\small
\renewcommand{\arraystretch}{1.2}
\begin{tabular}{
    p{0.49\textwidth}
    p{0.35\textwidth}
}
\hline
\textbf{Training component} &
\textbf{Configuration} \\
\hline

Training examples per run &
10,000 \\

Validation examples &
1,000 at the training length \\

Test examples &
1,000 at each evaluation length \\

Training epochs &
10 \\

Optimiser &
AdamW \\

Learning rate &
\(10^{-3}\) \\

Weight decay &
\(0.01\) \\

Learning-rate schedule &
None \\

Checkpoint selection &
Final training epoch \\

Addition batch size &
8 \\

Delayed Copy batch size &
1 \\

Reverse batch size &
1 \\

Random seeds &
42, 123, and 2024 \\

\hline
\end{tabular}
\caption{Training and evaluation configuration for the synthetic baseline
experiments.}
\label{tab:training-configuration}
\end{table}

The three seeds controlled Python, NumPy, and PyTorch randomness, including
synthetic data generation and model initialisation. GPU random-number
generators were also seeded when CUDA was available. The seed repetitions
therefore measured the combined sensitivity of the experimental pipeline to
data sampling, parameter initialisation, minibatch ordering, and stochastic
training operations.

After training, each checkpoint was loaded in evaluation mode and evaluated
separately on the held-out dataset for every configured test length. Dropout
was disabled during evaluation, and predictions were obtained using the
highest-logit token at each output position. All test sets were generated
before evaluation and were not used for model selection or further
optimisation.

\section{Evaluation Metrics}
\label{sec:evaluation-metrics}

The synthetic tasks were evaluated using token accuracy and exact-match
accuracy. Let \(y_{n,j}\) denote the target token at position \(j\) in
sequence \(n\), and let \(\hat{y}_{n,j}\) denote the corresponding predicted
token. A binary mask \(m_{n,j}\) indicates whether that position is included
in evaluation. For Addition and Delayed Copy, padding targets have
\(m_{n,j}=0\); for Reverse, every position has \(m_{n,j}=1\).

Token accuracy at test length \(L\) is defined as

\[
A_{\mathrm{tok}}(L)
=
100
\frac{
\sum_{n=1}^{N}
\sum_{j}
m_{n,j}
\mathbb{I}
\left[
\hat{y}_{n,j}=y_{n,j}
\right]
}{
\sum_{n=1}^{N}
\sum_{j}m_{n,j}
}.
\]

It measures the percentage of evaluated target tokens predicted correctly.
Padding positions do not contribute to either the numerator or denominator.

Exact-match accuracy is defined as

\[
A_{\mathrm{EM}}(L)
=
\frac{100}{N}
\sum_{n=1}^{N}
\mathbb{I}
\left[
\hat{y}_{n,j}=y_{n,j}
\text{ for every position satisfying }m_{n,j}=1
\right].
\]

A sequence is therefore counted as correct only when every non-padding target
token is predicted correctly. Exact-match accuracy was treated as the primary
synthetic-task metric because each task requires a complete output sequence.
Token accuracy was retained as a complementary metric because it distinguishes
complete sequence failure from outputs containing some correct tokens.

For interpretability, both metrics were also reported as error percentages:

\[
E_{\mathrm{tok}}(L)=100-A_{\mathrm{tok}}(L),
\qquad
E_{\mathrm{EM}}(L)=100-A_{\mathrm{EM}}(L).
\]

Token error was compared with a uniform-prediction reference over the active
target vocabulary. This reference is 90% error for Addition, which predicts
one of ten digits at active positions, and 99% error for Delayed Copy and
Reverse, which use 100 possible content symbols. These references apply to
token error and should not be confused with the 10% exact-match collapse
threshold defined below.

\subsection{Multi-seed aggregation}

For every task, positional encoding, and test length, performance was
aggregated across the three seeds. If the seed-level measurements are
\(A_1,A_2,A_3\), their mean is

\[
\bar{A}=\frac{1}{3}\sum_{s=1}^{3}A_s.
\]

The sample standard deviation was used to describe variation between seeds,
while the standard error was calculated as

\[
\mathrm{SE}=\frac{\mathrm{SD}}{\sqrt{3}}.
\]

Standard deviation and standard error were not treated as interchangeable.
Standard deviation describes the observed dispersion between the three runs,
whereas standard error describes uncertainty in their estimated mean. Every
result figure specifies explicitly which quantity, if any, is displayed by
its error bars or shaded region.

\subsection{Baseline eligibility}

Baseline eligibility was determined before extrapolation performance was
interpreted. A seed was considered to have fitted the training task when

\[
A_{\mathrm{EM}}(L_{\mathrm{train}})\geq 90\%.
\]

A task--encoding configuration was classified as \emph{eligible} only when
all three seeds met this threshold. A configuration was classified as
\emph{seed unstable} when the threshold was met by only one or two seeds.
Configurations for which no seed met the threshold were classified as
\emph{not eligible}. Only eligible task--encoding configurations were used
as the principal baselines for conclusions about extrapolation failure.

The 90\% criterion is an operational eligibility threshold rather than a claim
that 90\% constitutes perfect learning. Its purpose is to prevent
training-length underfit or seed instability from being misreported as
length-generalisation failure.

\section{Performance Degradation Curves}
\label{sec:performance-degradation}

Performance degradation curves plot accuracy or error percentage against the
nominal test length. They show the training-length result together with every
longer evaluation point, making it possible to distinguish an immediate
collapse at the first unseen length from degradation that emerges only at
larger lengths.

For each eligible task--encoding configuration, the first collapse length was
defined using the mean exact-match accuracy across seeds:

\[
L_{\mathrm{collapse}}
=
\min
\left\{
L>L_{\mathrm{train}}
:
\bar{A}_{\mathrm{EM}}(L)<10\%
\right\}.
\]

If mean exact-match accuracy fell below 10\% at the first unseen length, the
observed behaviour was classified as \emph{immediate extrapolation collapse}.
If the threshold was crossed only at a later length, the behaviour was
classified as delayed degradation followed by collapse. If no evaluated
length fell below the threshold, the result was described as \emph{no
observed collapse within the evaluated range}.

The 10\% value is an operational reporting threshold. It is not the theoretical
chance accuracy of an entire output sequence and does not, by itself,
identify the internal cause of failure. Its purpose is to provide a consistent
criterion for locating severe degradation across tasks whose complete output
lengths differ.

The eligible-baseline plots show the multi-seed mean together with one
standard deviation and faint individual-seed curves. This presentation makes
both the average length trend and seed-level stability visible. Training
length is marked using a vertical reference line, and the 10\% exact-match
threshold is marked using a horizontal reference line.

Because the evaluated lengths increase approximately geometrically, they are
displayed at equally spaced horizontal positions with their true numerical
values retained as tick labels. This improves label readability and prevents
the shorter lengths from being compressed near the left edge. The equal
visual spacing should therefore not be interpreted as a linear numerical
distance between consecutive sequence lengths.

\section{Attention Diagnostics}
\label{sec:attention-diagnostics}

Accuracy and error percentages identify performance changes but do not directly
describe how the model's attention distributions change with sequence length.
Attention statistics were therefore collected as supporting diagnostics for
the eligible baseline configurations. Because attention weights are not
necessarily faithful explanations of a model's prediction
~\cite{jain2019attention}, these statistics were not treated as a complete
mechanistic account of failure.

Attention was examined for the five eligible task--encoding configurations:
Addition with learned and sinusoidal encodings, Delayed Copy with sinusoidal
and RoPE encodings, and Reverse with learned positional embeddings. The
analysis used the seed-42 checkpoint as a representative run. For each
configuration, attention was extracted at the training length, the first
unseen length, and the longest evaluated length.

For an attention distribution \(\alpha_{ij}\), where \(i\) is the query
position and \(j\) is the key position, three statistics were calculated.
Normalised attention entropy was defined as

\[
H_{\mathrm{norm}}
=
\frac{
-\sum_{j=1}^{S}\alpha_{ij}\log(\alpha_{ij}+\epsilon)
}{
\log S
},
\]

where \(S\) is the encoded sequence length and
\(\epsilon=10^{-9}\) provides numerical stability. Values near zero indicate
concentrated attention, while values closer to one indicate a more diffuse
distribution. The statistic was averaged across query positions, attention
heads, examples, and Transformer layers.

Average attention distance was defined as

\[
D_{\mathrm{attn}}
=
\sum_{j=1}^{S}\alpha_{ij}|i-j|,
\]

and was averaged over query positions, heads, examples, and layers. This
measures the expected absolute positional separation between a query and the
positions receiving its attention.

Finally, the local attention ratio was defined as

\[
R_{\mathrm{local}}
=
\sum_{j:\,|i-j|\leq 10}\alpha_{ij}.
\]

It measures the proportion of attention assigned within a fixed window of ten
positions on either side of each query. The window was held constant at every
evaluation length.

Attention was collected from the first five test batches at each selected
length. Because the baseline batch size was eight for Addition and one for
Delayed Copy and Reverse, this corresponded to at most 40 Addition examples
and five examples for each of the other tasks. The attention analysis is
therefore a representative diagnostic rather than a full multi-seed
statistical evaluation.

Average distance and local-attention ratio also depend mechanically on
sequence length. A fixed local window represents a decreasing fraction of a
longer sequence, while an unnormalised distance has a larger possible range
at longer lengths. These measures were consequently interpreted as
descriptive changes accompanying failure, not as independent proof of its
cause.

\section{Mitigation Experiments}
\label{sec:mitigation-methodology}

Mitigation experiments were conducted after the baseline analysis. Each
intervention was associated with a specific hypothesis about the observed
failure behaviour rather than being treated as a general-purpose solution.
The experiments used seed 42 and were compared with the corresponding
seed-42 original baseline. They were therefore targeted follow-up experiments,
not a second complete multi-seed matrix.

Three principal interventions were examined: mixed-length training,
curriculum learning, and randomised padded training. Curriculum learning was
motivated by the principle of increasing task difficulty gradually during
optimisation~\cite{bengio2009curriculum}. Randomised padded training was
motivated by work examining whether variation in positional exposure can
improve length generalisation~\cite{ruoss2023randomized}. The implementation
used here was a narrower test and should not be interpreted as a direct
replication of that work.

\begin{table}[htbp]
\centering
\scriptsize
\renewcommand{\arraystretch}{1.2}
\begin{tabular}{
    p{0.19\textwidth}
    p{0.22\textwidth}
    p{0.27\textwidth}
    p{0.25\textwidth}
}
\hline
\textbf{Intervention} &
\textbf{Applied conditions} &
\textbf{Training design} &
\textbf{Hypothesis tested} \\
\hline

Mixed length &
Addition--Learned and Delayed Copy--Sinusoidal &
Joint exposure to lengths below and including the original training length &
Broader length exposure may reduce dependence on one fixed training length. \\

Curriculum &
Delayed Copy--Sinusoidal, Delayed Copy--RoPE, and Reverse--Learned &
Sequential stages at lengths 32, 64, 96, and 128 &
Gradual increases in length may overcome an optimisation barrier associated
with direct training at length 128. \\

Randomised padded &
Delayed Copy--Sinusoidal and Delayed Copy--RoPE &
Lengths 32, 64, 96, and 128 combined in one shuffled dataset and right-padded
to a common maximum &
Variation in separator and target positions may reduce reliance on their
fixed absolute locations. \\

Single-length control &
Delayed Copy--Sinusoidal &
Original length-128 distribution reproduced through the mitigation pipeline &
Any observed difference should result from the intervention rather than the
alternative training pipeline. \\

\hline
\end{tabular}
\caption{Mitigation conditions and the specific hypotheses they were designed
to test.}
\label{tab:mitigation-design}
\end{table}

For Addition, the first mixed-length variant distributed 10,000 training
examples across operand lengths 4, 8, 12, and 16 and trained for 10 epochs. A
strengthened variant used 40,000 examples and 20 epochs to test whether the
first condition was limited by data volume or optimisation time.

For Delayed Copy, mixed-length and randomised-padded conditions used lengths
32, 64, 96, and 128 with 40,000 examples in total. In the mixed-length
condition, each length was processed through a separate fixed-length data
loader. In the randomised-padded condition, all examples were combined,
shuffled, and right-padded to the maximum training sequence length. Padding
was not placed at random offsets; the intervention varied the source,
separator, and output locations through the sampled task length.

Curriculum models were trained sequentially at lengths 32, 64, 96, and 128.
Each stage used 10,000 examples and five epochs, and the final checkpoint from
one stage initialised the next. Only the final length-128 checkpoint was
evaluated over the complete test-length schedule.

The mitigation configurations used a batch size of eight. This differed from
the batch size of one used in the original Delayed Copy and Reverse baselines.
The single-length Delayed Copy control was therefore included to test whether
the modified training pipeline could reproduce the original sinusoidal
baseline without changing the length distribution.

Each intervention was first evaluated at the original training length. A
mitigation satisfying

\[
A_{\mathrm{EM}}(L_{\mathrm{train}})\geq 90\%
\]

was considered to have retained training fit and could therefore produce an
interpretable extrapolation comparison. If it fell below this threshold, its
longer-length outcome was classified as inconclusive due to training underfit.
For interventions that retained training fit, exact-match accuracy at the
first unseen and longest test lengths was compared directly with the
corresponding original baseline. The phrase \emph{no observed improvement}
is used when an eligible intervention did not increase extrapolation
performance within the evaluated range.

The Reverse curriculum configuration used a 102-symbol vocabulary, whereas
the original Reverse baseline used 100 symbols. Unless this condition is
rerun using the baseline vocabulary, it is not an exactly matched
baseline-versus-intervention comparison and is treated only as an exploratory
curriculum sensitivity test.

\section{Real-World Extension}
\label{sec:real-world-extension}

The synthetic benchmark provides controlled tests in which the target rule is
known exactly. A separate real-world experiment was included to examine
whether input-length sensitivity also appears in multivariate time-series
classification. It was treated as an exploratory contrast rather than
validation of the synthetic failure taxonomy.

The experiment used SelfRegulationSCP2 from the UEA Multivariate Time Series
Classification Archive~\cite{bagnall2018uea}. The dataset contains EEG
recordings from a binary cursor-control task. It provides 200 training
examples and 180 test examples, with seven channels and 1152 time points per
example. The predefined archive split was retained.

Each channel was standardised using its mean and standard deviation calculated
from the training split across examples and time points. The same
training-derived statistics were then applied to the test split. No
information from the test set was used to calculate the normalisation
parameters.

The classifier was trained using the first 256 time points of every training
series. Evaluation used prefixes of lengths 256, 512, 1024, and 1152 from
the predefined test examples. Unlike the synthetic protocol, increasing the
prefix length introduces additional observed measurements and may therefore
add useful information or noise. The experiment measures sensitivity to
prefix length rather than deterministic rule extrapolation.

The time-series model replaced the synthetic token-embedding layer with a
linear projection from the seven input channels to the 128-dimensional model
representation. It then applied the same two Transformer blocks and positional
conditions used in the synthetic experiments. The representations were
mean-pooled across time and passed to a linear two-class classifier.

Models were trained for 50 epochs using AdamW, a learning rate of
\(10^{-3}\), batch size 16, and cross-entropy loss. The five positional
conditions were repeated with seeds 42, 123, and 2024, producing 15 trained
classifiers. Classification accuracy was calculated on the 180-example test
split at each prefix length. Mean accuracy and standard error were then
calculated across the three seeds.

\begin{table}[htbp]
\centering
\small
\renewcommand{\arraystretch}{1.2}
\begin{tabular}{
    p{0.43\textwidth}
    p{0.42\textwidth}
}
\hline
\textbf{Component} &
\textbf{Setting} \\
\hline

Dataset &
SelfRegulationSCP2 \\

Data type &
Seven-channel EEG time series \\

Classes &
2 \\

Training examples &
200 \\

Test examples &
180 \\

Full sequence length &
1152 time points \\

Training prefix &
256 time points \\

Evaluation prefixes &
256, 512, 1024, and 1152 \\

Input projection &
7 channels to 128 dimensions \\

Sequence aggregation &
Mean pooling across time \\

Training epochs &
50 \\

Batch size &
16 \\

Learning rate &
\(10^{-3}\) \\

Positional conditions &
Learned, sinusoidal, NoPE, ALiBi, and RoPE \\

Random seeds &
42, 123, and 2024 \\

Evaluation metric &
Classification accuracy \\

\hline
\end{tabular}
\caption{Configuration of the exploratory SelfRegulationSCP2 experiment.}
\label{tab:real-world-setup}
\end{table}

The synthetic 90% baseline-eligibility and 10% exact-match collapse thresholds
were not applied to this experiment because its output is a single noisy
class label rather than a deterministic token sequence. Its results were
therefore interpreted separately and were not used to determine whether a
synthetic failure category had been reproduced.

\section{Failure Taxonomy Construction}
\label{sec:failure-taxonomy-construction}

The final methodological stage organised the observed behaviours into a
descriptive failure taxonomy. The taxonomy was constructed sequentially from
training-length fit, seed stability, extrapolation curves, token-error
percentages, position-wise error profiles, task-specific diagnostics, and
mitigation comparisons. Attention statistics were retained as secondary
context but were not used alone to assign a failure category.

The construction followed five questions:

\begin{enumerate}
    \item Did the configuration learn the training-length task?
    \item Was training-length performance stable across all three seeds?
    \item If the baseline was eligible, where did severe extrapolation
    degradation first occur?
    \item How were token errors distributed across output positions and
    task-specific conditions?
    \item Did an intervention retain training fit and improve performance at
    unseen lengths relative to its corresponding baseline?
\end{enumerate}

These questions produced the five behavioural categories summarised in
Table~\ref{tab:failure-taxonomy-method}. The categories are not all mutually
exclusive. For example, an eligible configuration may exhibit immediate
extrapolation collapse, have an approximately uniform position-wise error
profile, and show no observed improvement under a tested intervention.

\begin{table}[htbp]
\centering
\scriptsize
\renewcommand{\arraystretch}{1.2}
\begin{tabular}{
    p{0.22\textwidth}
    p{0.34\textwidth}
    p{0.35\textwidth}
}
\hline
\textbf{Behavioural category} &
\textbf{Operational criterion} &
\textbf{Interpretive role} \\
\hline

Training-length underfitting &
The configuration does not reach 90\% exact match at the training length &
Prevents failure to learn the original task from being misreported as
extrapolation failure. \\

Seed instability &
The 90\% training-fit threshold is met by some, but not all, seeds &
Identifies conclusions that depend on data sampling, initialisation, or
stochastic optimisation. \\

Immediate extrapolation collapse &
An eligible configuration has mean exact-match accuracy below 10\% at the
first unseen length &
Locates severe degradation directly beyond the observed training-length
boundary without assigning a specific internal cause. \\

Spatial error structure &
Token errors are described by their distribution across relative output
positions, carry conditions, or dependency distances &
Distinguishes approximately uniform near-chance collapse from more
position-dependent failure behaviour. \\

No observed improvement under tested mitigation &
An intervention retains at least 90\% training-length exact match but does not
increase performance at the evaluated unseen lengths relative to its baseline &
Reports the outcome of the specific intervention without claiming that the
failure is generally resistant to mitigation. \\

\hline
\end{tabular}
\caption{Operational construction of the failure taxonomy.}
\label{tab:failure-taxonomy-method}
\end{table}

Mitigation conditions that failed to retain training fit were classified as
inconclusive due to underfitting rather than as unsuccessful evidence about
extrapolation. Similarly, changes in attention entropy, distance, or locality
were treated only as accompanying descriptive observations. This prevented
the taxonomy from making stronger mechanistic claims than were supported by
the experimental evidence.

The taxonomy is therefore an evidence-organising framework rather than a
statistical proof that each category corresponds to a unique internal model
mechanism. Its purpose is to separate observably different outcomes that
would otherwise be obscured by reporting only accuracy at the longest test
length.

\section{Output Organisation}
\label{sec:output-organisation}

Experimental artefacts were separated into configurations, checkpoints, logs,
raw result files, derived analyses, and report-ready figures. File names
encode the task or dataset, training length, positional condition, and random
seed where applicable. This makes individual results traceable to the
configuration and checkpoint that produced them.

Raw synthetic baseline results are stored in
\texttt{outputs/results/multiseed\_results/}, mitigation results in
\texttt{outputs/results/mitigation\_results/}, and time-series results in
\texttt{outputs/results/realworld/}. These CSV files contain the unaggregated
measurements for individual task--encoding--seed combinations.

Derived attention statistics are stored in \texttt{outputs/analysis/attention/} and \texttt{outputs/analysis/atte-} \\ \texttt{ntion\_summary/}. Position-wise diagnostic CSV files transferred from the GPU environment are retained under \texttt{position\_diagnostics\_aws/final/}. Mitigation and taxonomy summaries
are stored under \texttt{outputs/analysis/}.

Report-ready PDF and PNG figures, together with the summary CSV files used to
construct them, are stored in \texttt{outputs/plots/final/}. The PDF versions
were used in the dissertation to preserve vector text and line rendering,
while PNG versions were retained for rapid inspection.

\begin{table}[htbp]
\centering
\small
\renewcommand{\arraystretch}{1.2}
\begin{tabular}{
    p{0.29\textwidth}
    p{0.54\textwidth}
}
\hline
\textbf{Artefact} &
\textbf{Location} \\
\hline

Experiment configurations &
\texttt{configs/multiseed/},
\texttt{configs/mitigation/}, and
\texttt{configs/realworld/} \\

Model checkpoints &
\texttt{outputs/checkpoints/} \\

Execution logs &
\texttt{outputs/logs/} \\

Raw baseline results &
\texttt{outputs/results/multiseed\_results/} \\

Raw mitigation results &
\texttt{outputs/results/mitigation\_results/} \\

Raw real-world results &
\texttt{outputs/results/realworld/} \\

Attention diagnostics &
\texttt{outputs/analysis/attention/} and
\texttt{outputs/analysis/attention\_summary/} \\

Position-wise diagnostics &
\texttt{position\_diagnostics\_aws/final/} \\

Final figures and report summaries &
\texttt{outputs/plots/final/} \\

\hline
\end{tabular}
\caption{Organisation of experimental artefacts and derived report outputs.}
\label{tab:output-organisation}
\end{table}

Raw CSV files were retained separately from aggregated summaries so that
reported means, error percentages, eligibility decisions, and plots could be
checked against their underlying seed-level measurements. Plotting and
analysis scripts read these stored results rather than manually transcribed
values.

\section{Implementation and Compute Setup}
\label{sec:implementation-compute}

The experimental pipeline was implemented in Python using PyTorch
~\cite{paszke2019pytorch}. NumPy was used for numerical preprocessing, pandas
for result aggregation, Matplotlib for figure generation, PyYAML for
configuration loading, and aeon for access to the UEA time-series dataset.

Synthetic task generators were implemented as separate Python classes.
Generated examples were stored in JSON Lines format, with one input--target
pair per record. Experimental settings were stored in YAML configuration
files rather than embedded directly in training commands. Training,
evaluation, diagnostic analysis, and figure generation were implemented as
separate executable modules.

\begin{table}[htbp]
\centering
\small
\renewcommand{\arraystretch}{1.2}
\begin{tabular}{
    p{0.42\textwidth}
    p{0.43\textwidth}
}
\hline
\textbf{Project component} &
\textbf{Repository location} \\
\hline

Synthetic task generation &
\texttt{src/tasks/} \\

Dataset loading and preprocessing &
\texttt{src/data/} \\

Transformer and positional implementations &
\texttt{src/models/} \\

Training, evaluation, and analysis scripts &
\texttt{src/training/} \\

Baseline configurations &
\texttt{configs/multiseed/} \\

Mitigation configurations &
\texttt{configs/mitigation/} \\

Real-world configurations &
\texttt{configs/realworld/} \\

Dependency specification &
\texttt{requirements.txt} \\

Project instructions and commands &
\texttt{README.md} \\

\hline
\end{tabular}
\caption{Source-code and configuration components supporting reproducibility.}
\label{tab:reproducibility-components}
\end{table}

Local macOS execution was used for implementation, debugging, result
aggregation, and figure generation. GPU-intensive training and diagnostic
evaluation were performed on an Amazon Web Services G5-family instance using
CUDA. The same Git repository and YAML-controlled experiment interface were
used in both environments.

Random seeds were recorded in every experimental configuration, while
training and evaluation logs preserved the executed condition and output
location. Checkpoints and raw CSV results were transferred from the GPU
environment before the final summaries and figures were generated locally.

The repository targets Python 3.11 and records its required packages in
\texttt{requirements.txt}. Because that file lists dependency names without
exact version pins, complete byte-for-byte environment reconstruction is not
guaranteed. This is a reproducibility limitation; the configurations, random
seeds, raw results, checkpoints, and analysis scripts nevertheless preserve
the principal experimental workflow and reported measurements.


% -----------------------------------------------------------------------------

\chapter{Synthetic Benchmark Results}
\label{chap:synthetic-results}

\noindent
This chapter reports the results of the controlled synthetic benchmark. The
analysis first establishes which task--encoding configurations learned the
training-length task reliably across all three seeds. Extrapolation behaviour
is then reported only for these eligible baselines, while configurations that
underfit or exhibited seed instability are analysed separately.

Five of the fifteen task--encoding configurations satisfied the multi-seed
baseline-eligibility criterion. Every one of these eligible baselines fell
below the 10\% exact-match threshold at the first unseen length. The principal
result is therefore not simply that some models performed poorly at long
lengths, but that configurations with stable training-length fit underwent
immediate and reproducible extrapolation collapse when the nominal task length
first exceeded the training boundary.

The chapter first reports training fit and seed stability across the complete
baseline matrix. It then examines eligible-baseline extrapolation and error
percentages for Addition, Delayed Copy, and Reverse. Position-wise and
task-specific error structure is analysed separately in
Chapter~\ref{chap:failure-analysis}.


\section{Overview of Synthetic Results}
\label{sec:synthetic-results-overview}

The baseline matrix contained fifteen task--encoding configurations, each
evaluated using seeds 42, 123, and 2024. A task--encoding configuration was
eligible for the principal extrapolation analysis only when all three seeds
achieved at least 90\% exact-match accuracy at the training length.

Figure~\ref{fig:baseline-training-fit} shows the training-length exact-match
accuracy of every seed. The point shape identifies the seed, the horizontal
black segment gives the mean, and the dashed line marks the 90\% eligibility
threshold.

\begin{figure}[htbp]
    \centering
    \includegraphics[width=\textwidth]
    {baseline_training_fit_by_seed.pdf}
    \caption{Training-length exact-match accuracy across tasks, positional
    encodings, and seeds. Circles, squares, and triangles represent seeds 42,
    123, and 2024 respectively; horizontal black segments show the
    three-seed mean. The dashed line marks the 90\% baseline-eligibility
    threshold, and annotations report the number of seeds satisfying it.}
    \label{fig:baseline-training-fit}
\end{figure}

Five configurations provided stable eligible baselines: learned and
sinusoidal encodings for Addition, sinusoidal and RoPE encodings for Delayed
Copy, and learned positional embeddings for Reverse. Delayed Copy with
learned embeddings was seed unstable because only seed 123 learned the
training task. Reverse with sinusoidal encoding was also seed unstable because
seeds 42 and 123 succeeded while seed 2024 failed.

The remaining configurations did not satisfy the eligibility criterion.
NoPE and ALiBi failed to learn the training task across all three tasks.
Addition with RoPE achieved partial training fit but remained below the
eligibility threshold for every seed, while Reverse with RoPE also failed at
the training length. These configurations cannot support conclusions about
length extrapolation because the original task was not first learned
reliably.

The eligible configurations showed a common sequence-level outcome: mean
exact-match accuracy fell to 0\% at the first unseen length and remained at
0\% throughout the evaluated range. Token-level results, however, differed
between Addition and the two 100-symbol tasks. The following sections
therefore report both sequence and token error rather than describing every
0\% exact-match result as the same form of failure.


\section{Addition Results}
\label{sec:addition-results}

Addition models were trained using operands of length 16 and evaluated using
operand lengths from 16 to 1024. Learned and sinusoidal positional encodings
were the only Addition configurations that satisfied the multi-seed
eligibility criterion.

At the training length, learned positional embeddings achieved a mean
exact-match accuracy of 92.0\%, with a standard deviation of 1.84 percentage
points and a seed range of 90.7--94.1\%. Sinusoidal encoding achieved a mean
of 93.7\%, with a standard deviation of 0.95 percentage points and a range of
92.8--94.7\%. All six corresponding seed-level runs therefore exceeded the
90\% eligibility threshold.

Figure~\ref{fig:addition-eligible-extrapolation} shows the extrapolation curves
for these two eligible baselines. Both fell to 0\% exact-match accuracy at
length 32, the first unseen length, and remained at 0\% through length 1024.
The collapse was reproduced across all three seeds rather than being caused by
one unstable run.

\begin{figure}[htbp]
    \centering
    \includegraphics[width=\textwidth]
    {addition_eligible_baseline_extrapolation.pdf}
    \caption{Exact-match accuracy of the eligible Addition baselines across
    operand lengths. Bold curves show the three-seed mean, variability is
    shown as mean \(\pm\) one standard deviation, and faint curves show
    individual seeds. The vertical dashed line marks the training length and
    the horizontal dotted line marks the 10\% operational collapse
    threshold.}
    \label{fig:addition-eligible-extrapolation}
\end{figure}

This result is classified as immediate extrapolation collapse: both
configurations learned the length-16 task consistently, but their mean
exact-match accuracy crossed the collapse threshold at the first unseen
length. The accuracy curve establishes the location and severity of the
failure, but it does not by itself identify the internal strategy learned by
the models.

The non-eligible Addition conditions exhibited different training-length
behaviour. NoPE and ALiBi achieved 0\% mean exact-match accuracy at length 16.
RoPE achieved a mean of 47.3\%, with a standard deviation of 3.60 percentage
points and a range of 44.1--51.2\%. RoPE therefore demonstrated partial
training fit but did not provide a valid multi-seed extrapolation baseline.

Exact-match accuracy alone does not show whether the failed outputs contained
mostly correct digits or were incorrect throughout. Figure
~\ref{fig:addition-error-decomposition} therefore presents sequence and token
error percentages for the eligible configurations.

\begin{figure}[htbp]
    \centering
    \includegraphics[width=\textwidth]
    {addition_eligible_baseline_error_decomposition.pdf}
    \caption{Sequence-level and token-level error for the eligible Addition
    baselines. Sequence error is \(100-A_{\mathrm{EM}}\), while token error is
    \(100-A_{\mathrm{tok}}\). Bold curves show the three-seed mean, variability
    is shown as mean \(\pm\) one standard deviation, and faint curves show
    individual seeds. The token-error panel includes the 90\% error reference
    corresponding to uniform prediction over the ten active digit symbols.}
    \label{fig:addition-error-decomposition}
\end{figure}

At the training length, mean sequence error was 8.0\% for learned embeddings
and 6.3\% for sinusoidal encoding. Their corresponding token errors were only
0.52\% and 0.40\%, showing that the remaining sequence errors were associated
with a small proportion of incorrect digits.

At length 32, sequence error increased to 100\% for both eligible
configurations. Token error increased to 87.74\% for learned embeddings and
87.16\% for sinusoidal encoding. Token error then approached the 90\%
uniform-prediction reference as length increased, reaching 89.93\% and
89.98\% respectively at length 1024.

Therefore, the 0\% exact-match result was not produced by otherwise
near-correct sums containing only one occasional error. At unseen lengths,
the majority of active output digits were incorrect, and performance moved
towards the uniform digit-prediction reference. The spatial distribution of
these errors and their relationship to incoming carry are examined in
Chapter~\ref{chap:failure-analysis}.

\section{Delayed Copy Results}
\label{sec:copy-results}

As specified in Section~\ref{sec:method-synthetic-tasks}, Delayed Copy
requires the model to reproduce each content token after a separator and a
delay. Models were trained at a nominal length of 128 and evaluated at lengths
256, 512, and 1024. The encoded sequence length was \(2L+1\), so increasing
the nominal length also increased the positional separation between each
source token and its corresponding target position.

Sinusoidal encoding and RoPE were the two eligible Delayed Copy baselines.
Sinusoidal encoding achieved a mean training-length exact-match accuracy of
99.93\%, with a standard deviation of 0.12 percentage points and a seed range
of 99.8--100.0\%. RoPE achieved 100\% exact-match accuracy for all three
seeds.

Figure~\ref{fig:copy-eligible-extrapolation} shows that both configurations
fell to 0\% exact-match accuracy at length 256, the first unseen length. They
remained at 0\% at lengths 512 and 1024. The result was reproduced across all
three seeds for both positional encodings.

\begin{figure}[htbp]
    \centering
    \includegraphics[width=\textwidth]
    {copy_eligible_baseline_extrapolation.pdf}
    \caption{Exact-match accuracy of the eligible Delayed Copy baselines
    across nominal sequence lengths. Bold curves show the three-seed mean,
    variability is shown as mean \(\pm\) one standard deviation, and faint
    curves show individual seeds. The vertical dashed line marks the training
    length and the horizontal dotted line marks the 10\% operational collapse
    threshold.}
    \label{fig:copy-eligible-extrapolation}
\end{figure}

\begin{figure}[htbp]
    \centering
    \includegraphics[width=\textwidth]
    {copy_eligible_baseline_error_decomposition.pdf}
    \caption{Sequence-level and token-level error for the eligible Delayed
    Copy baselines. Bold curves show the three-seed mean, variability is
    shown as mean \(\pm\) one standard deviation, and faint curves show
    individual seeds. The token-error panel includes the 99\% error reference
    corresponding to uniform prediction over the 100 content symbols.}
    \label{fig:copy-error-decomposition}
\end{figure}

This behaviour is classified as immediate extrapolation collapse. The models
reproduced almost every training-length sequence exactly, but neither
configuration produced one completely correct sequence among the 3,000
seed-level test examples at length 256. The result locates the failure at the
first unseen length without requiring an assumption about the internal
retrieval strategy used by the model.

The token-level results show that the failed sequences were not close to
correct. Figure~\ref{fig:copy-error-decomposition} presents sequence and token
error percentages for the eligible configurations.



At the training length, mean sequence error was 0.07\% for sinusoidal
encoding and 0\% for RoPE. Their token errors were approximately 0.0005\% and
0\%, respectively. At length 256, sequence error increased to 100\% for both
configurations, while token error increased to 99.00\% for sinusoidal encoding
and 98.98\% for RoPE. Token errors remained close to the 99\% uniform-symbol
reference at lengths 512 and 1024.

Therefore, the failure was not limited to a small number of misplaced tokens.
At unseen lengths, predictions across the active target positions were
approximately as inaccurate as uniform prediction over the content
vocabulary.

Learned positional embeddings did not provide an eligible pooled baseline.
Only seed 123 learned the training task, achieving 99.4\% exact-match accuracy,
while seeds 42 and 2024 achieved 0\%. This produced a mean of 33.13\% and a
standard deviation of 57.39 percentage points. The configuration is therefore
classified as seed unstable, and its pooled mean is not interpreted as a
representative extrapolation curve.

NoPE and ALiBi achieved 0\% exact-match accuracy at the training length for
all three seeds. Their longer-length results are consequently classified as
not applicable to extrapolation rather than as additional evidence of
length-generalisation failure.

\section{Reverse Results}
\label{sec:reverse-results}

Reverse requires each output position to retrieve the input token measured
from the opposite end of the sequence. Models were trained at length 128 and
evaluated at lengths 256, 512, and 1024. Learned positional embeddings were
the only Reverse configuration that satisfied the multi-seed eligibility
criterion.

The learned condition achieved 100\% exact-match and token accuracy at the
training length for all three seeds. Figure
~\ref{fig:reverse-eligible-extrapolation} shows that mean exact-match accuracy
then fell to 0\% at length 256 and remained at 0\% through length 1024. No
seed produced a non-zero exact-match result at any unseen length.

\begin{figure}[htbp]
    \centering
    \includegraphics[width=\textwidth]
    {reverse_eligible_baseline_extrapolation.pdf}
    \caption{Exact-match accuracy of the eligible Reverse baseline across
    sequence lengths. The bold curve shows the three-seed mean, variability
    is shown as mean \(\pm\) one standard deviation, and faint curves show
    individual seeds. The vertical dashed line marks the training length and
    the horizontal dotted line marks the 10\% operational collapse
    threshold.}
    \label{fig:reverse-eligible-extrapolation}
\end{figure}

\begin{figure}[htbp]
    \centering
    \includegraphics[width=\textwidth]
    {reverse_eligible_baseline_error_decomposition.pdf}
    \caption{Sequence-level and token-level error for the eligible Reverse
    baseline. The bold curve shows the three-seed mean, variability is shown
    as mean \(\pm\) one standard deviation, and faint curves show individual
    seeds. The token-error panel includes the 99\% error reference
    corresponding to uniform prediction over the 100-symbol vocabulary.}
    \label{fig:reverse-error-decomposition}
\end{figure}

This is another case of immediate extrapolation collapse. Unlike a
training-underfit configuration, the learned baseline solved every
training-length test example across all seeds before failing at the first
unseen length. The result demonstrates that stable in-distribution performance
did not extend across the change from length 128 to 256.

Figure~\ref{fig:reverse-error-decomposition} shows that the sequence-level
collapse was accompanied by token errors close to the uniform-symbol
reference.

Mean token error increased from 0\% at the training length to 98.75\% at
length 256. It remained close to the 99\% uniform-symbol reference, reaching
98.84\% at length 512 and 98.92\% at length 1024. The model therefore did not
merely produce reversals containing a small number of local mistakes; most
output positions were incorrect once the sequence length exceeded the
training boundary.

Sinusoidal encoding was seed unstable. Seeds 42 and 123 achieved 100\%
training-length exact-match accuracy, while seed 2024 achieved 0\%. This
produced a pooled mean of 66.6\% and a standard deviation of 57.68 percentage
points. Pooling these qualitatively different runs would obscure the
instability, so sinusoidal encoding was not treated as an eligible
three-seed baseline.

NoPE, ALiBi, and RoPE achieved 0\% training-length exact-match accuracy for
all three seeds. Their results are classified as training-length underfitting
and do not support conclusions about extrapolation. The distribution of
Reverse errors by relative output position and dependency distance is examined
in Chapter~\ref{chap:failure-analysis}.

\newpage
\section{Positional Encoding Comparison}
\label{sec:positional-comparison}

The positional-encoding conditions produced substantially different
training-length outcomes. However, their effects were task dependent rather
than forming a single consistent ranking. Learned positional embeddings were
eligible for Addition and Reverse but seed unstable for Delayed Copy.
Sinusoidal encoding was eligible for Addition and Delayed Copy but seed
unstable for Reverse. RoPE was eligible only for Delayed Copy.

Table~\ref{tab:positional-comparison} reports the training-length result for
all fifteen task--encoding configurations. The number of successful seeds
shows how many runs reached the 90\% exact-match eligibility threshold.
Training accuracy is reported as the three-seed mean and sample standard
deviation.

\begin{table}[htbp]
\centering
\scriptsize
\renewcommand{\arraystretch}{1.16}
\begin{tabular}{
    p{0.17\textwidth}
    p{0.14\textwidth}
    p{0.10\textwidth}
    p{0.18\textwidth}
    p{0.32\textwidth}
}
\hline
\textbf{Task} &
\textbf{Encoding} &
\textbf{Fit seeds} &
\textbf{Exact at train (\%)} &
\textbf{Baseline interpretation} \\
\hline

Addition &
Learned &
3/3 &
\(92.00 \pm 1.84\) &
Eligible; immediate collapse \\

Addition &
Sinusoidal &
3/3 &
\(93.70 \pm 0.95\) &
Eligible; immediate collapse \\

Addition &
NoPE &
0/3 &
\(0.00 \pm 0.00\) &
Training-length underfit \\

Addition &
ALiBi &
0/3 &
\(0.00 \pm 0.00\) &
Training-length underfit \\

Addition &
RoPE &
0/3 &
\(47.30 \pm 3.60\) &
Partial fit; not eligible \\
\hline

Delayed Copy &
Learned &
1/3 &
\(33.13 \pm 57.39\) &
Seed unstable \\

Delayed Copy &
Sinusoidal &
3/3 &
\(99.93 \pm 0.12\) &
Eligible; immediate collapse \\

Delayed Copy &
NoPE &
0/3 &
\(0.00 \pm 0.00\) &
Training-length underfit \\

Delayed Copy &
ALiBi &
0/3 &
\(0.00 \pm 0.00\) &
Training-length underfit \\

Delayed Copy &
RoPE &
3/3 &
\(100.00 \pm 0.00\) &
Eligible; immediate collapse \\
\hline

Reverse &
Learned &
3/3 &
\(100.00 \pm 0.00\) &
Eligible; immediate collapse \\

Reverse &
Sinusoidal &
2/3 &
\(66.60 \pm 57.68\) &
Seed unstable \\

Reverse &
NoPE &
0/3 &
\(0.00 \pm 0.00\) &
Training-length underfit \\

Reverse &
ALiBi &
0/3 &
\(0.00 \pm 0.00\) &
Training-length underfit \\

Reverse &
RoPE &
0/3 &
\(0.00 \pm 0.00\) &
Training-length underfit \\

\hline
\end{tabular}
\caption{Training-length performance and eligibility across the complete
synthetic baseline matrix. Exact-match values are the mean \(\pm\) one sample
standard deviation over three seeds. A fit seed achieved at least 90\%
training-length exact-match accuracy.}
\label{tab:positional-comparison}
\end{table}

Five configurations were eligible, two were seed unstable, and eight were not
eligible. All five eligible configurations subsequently achieved 0\% mean
exact-match accuracy at the first unseen length and at the longest evaluated
length. Therefore, none of the positional conditions produced reliable
sequence-level extrapolation within an eligible task--encoding configuration.

The non-eligible results must be interpreted differently. NoPE and ALiBi
failed to meet the training-fit criterion on every task in this particular
architecture and optimisation setting. This finding does not establish that
these methods generally cannot support length extrapolation; it shows that
their present configurations did not provide valid trained baselines from
which extrapolation could be assessed.

RoPE also displayed a strongly task-dependent outcome. It achieved 100\%
training-length exact match across all seeds on Delayed Copy, partial
training-length performance on Addition, and 0\% exact match on Reverse. The
comparison therefore provides no evidence that one positional encoding is
uniformly preferable across the three task structures.

Finally, although every configuration had 0\% mean exact-match accuracy at
the longest test length, this common endpoint does not imply a common failure
type. For eligible configurations, it represents extrapolation collapse. For
seed-unstable configurations, pooled results conceal different training
outcomes. For non-eligible configurations, it primarily reflects failure to
establish the training task.

\section{Failure Length Analysis}
\label{sec:failure-length-analysis}

The first collapse length was calculated only for configurations that
satisfied the multi-seed baseline-eligibility criterion. As defined in
Section~\ref{sec:performance-degradation}, it is the first evaluated unseen
length at which mean exact-match accuracy falls below 10\%.

Table~\ref{tab:first-collapse-lengths} reports the result for all five eligible
baselines. Addition crossed the threshold at length 32, while Delayed Copy
and Reverse crossed it at length 256. In every case, this was the first
evaluated length beyond \(L_{\mathrm{train}}\).

\begin{table}[htbp]
\centering
\small
\renewcommand{\arraystretch}{1.2}
\begin{tabular}{
    p{0.23\textwidth}
    p{0.20\textwidth}
    p{0.15\textwidth}
    p{0.17\textwidth}
    p{0.18\textwidth}
}
\hline
\textbf{Task} &
\textbf{Encoding} &
\(\boldsymbol{L_{\mathrm{train}}}\) &
\textbf{First unseen length} &
\textbf{First collapse length} \\
\hline

Addition &
Learned &
16 &
32 &
32 \\

Addition &
Sinusoidal &
16 &
32 &
32 \\

Delayed Copy &
Sinusoidal &
128 &
256 &
256 \\

Delayed Copy &
RoPE &
128 &
256 &
256 \\

Reverse &
Learned &
128 &
256 &
256 \\

\hline
\end{tabular}
\caption{First collapse length for every eligible synthetic baseline. Collapse
is defined as mean exact-match accuracy below 10\% at an unseen length.}
\label{tab:first-collapse-lengths}
\end{table}

All five eligible baselines therefore exhibited immediate extrapolation
collapse according to the operational definition. No eligible configuration
showed a gradual sequence-level decline across multiple unseen evaluation
points before crossing the threshold. Instead, exact-match accuracy changed
from at least 92\% at the training length to 0\% at the first unseen length.

The term \emph{immediate} refers to the first sampled unseen length rather than
to every possible length directly above the training boundary. Addition was
not evaluated at operand lengths between 16 and 32, and Delayed Copy and
Reverse were not evaluated at lengths between 128 and 256. The experiment
therefore establishes that collapse had occurred by the first evaluated
doubling of length, but it cannot identify the precise intermediate length at
which the degradation began.

Addition with RoPE is not assigned a collapse length. Although its mean
training-length exact-match accuracy was 47.3\%, none of its seeds reached the
90\% eligibility threshold. Labelling its length-32 result as extrapolation
collapse would incorrectly combine partial training fit with failure beyond
the training distribution.

The common collapse boundary shows that stable training-length performance did
not persist to the first sampled unseen length. It does not, by itself,
establish that the models learned representations explicitly tied to one
length; determining the internal cause would require stronger mechanistic
evidence.

\section{Multi-Seed Robustness}
\label{sec:multi-seed-robustness}

The multi-seed experiment examined sensitivity to synthetic data sampling,
parameter initialisation, minibatch ordering, and stochastic optimisation.
Figure~\ref{fig:baseline-training-fit}, presented in the chapter overview,
shows the complete seed-level training results. It should not be replaced or
repeated by a second mean-and-standard-error figure.

The five eligible configurations were stable across seeds: each of their
individual runs exceeded the 90\% training-fit threshold. However, Delayed
Copy with learned positional embeddings and Reverse with sinusoidal encoding
produced qualitatively different outcomes across seeds, as shown in
Table~\ref{tab:seed-unstable-configurations}.

\begin{table}[htbp]
\centering
\small
\renewcommand{\arraystretch}{1.0}
\begin{tabular}{
    p{0.22\textwidth}
    p{0.19\textwidth}
    p{0.13\textwidth}
    p{0.13\textwidth}
    p{0.13\textwidth}
    p{0.13\textwidth}
}
\hline
\textbf{Task} &
\textbf{Encoding} &
\textbf{Seed 42} &
\textbf{Seed 123} &
\textbf{Seed 2024} &
\textbf{Fit seeds} \\
\hline

Delayed Copy &
Learned &
0.0\% &
99.4\% &
0.0\% &
1/3 \\

Reverse &
Sinusoidal &
100.0\% &
99.8\% &
0.0\% &
2/3 \\

\hline
\end{tabular}
\caption{Training-length exact-match accuracy for the two seed-unstable
configurations. A fit seed achieved at least 90\% exact-match accuracy.}
\label{tab:seed-unstable-configurations}
\end{table}

For Delayed Copy with learned embeddings, only seed 123 learned the
training-length task. Pooling the three runs produced a mean of 33.13\% and a
standard deviation of 57.39 percentage points, but the mean does not describe
a typical run: it combines one nearly perfect model with two models at 0\%.

The same issue occurred for Reverse with sinusoidal encoding. Seeds 42 and
123 achieved 100.0\% and 99.8\%, respectively, while seed 2024 achieved
0\%. The resulting mean of 66.6\% and standard deviation of 57.68 percentage
points again conceal two distinct outcomes.

These results justify reporting individual seeds alongside aggregated
statistics. Mean performance alone would make both configurations appear to
have achieved partial training fit, when the seed-level evidence instead
shows successful and unsuccessful optimisation outcomes. For this reason,
neither configuration was pooled as an eligible baseline.

Seed stability must also be separated from task success. NoPE and ALiBi were
stable in the sense that their training-length exact-match accuracy was
consistently 0\%, but this consistency does not make them valid extrapolation
baselines. Robustness therefore required both adequate training performance
and agreement across seeds.

\section{Generalisation Gap}
\label{sec:generalisation-gap}
For descriptive purposes, the training-to-longest generalisation gap can be
defined as $\bar{A}{\mathrm{EM}}(L{\mathrm{max}})$.

A large positive value can indicate that performance was high at the training
length but low at the longest evaluated length. However, the gap is only
meaningful after training fit has been established. A model scoring 0% at
both lengths has a gap of zero, but this reflects training underfit rather
than successful generalisation.
In the present benchmark, mean exact-match accuracy at the longest length was
0% for all fifteen task--encoding configurations. Consequently, each
configuration's numerical gap was identical to its training-length mean. The
gap therefore added no information beyond the baseline-eligibility analysis
and was not used as a primary ranking statistic.
Among the eligible baselines, the gaps were 92.0 percentage points for
Addition with learned embeddings, 93.7 for Addition with sinusoidal encoding,
99.93 for Delayed Copy with sinusoidal encoding, and 100 for both Delayed Copy
with RoPE and Reverse with learned embeddings. These values quantify the
difference between the endpoints, but they do not show where failure began.
The first-collapse analysis is more informative because it establishes that
all five eligible configurations crossed the operational threshold at the
first sampled unseen length.
The legacy generalisation-gap plot is therefore omitted from the final report.
Baseline eligibility, seed-level training fit, and complete degradation curves
provide a clearer and less potentially misleading account of the results.

\section{Comparison with Expected Positional Encoding Behaviour}
\label{sec:expected-positional-behaviour}

The results do not support a simple expectation that selecting a positional
encoding associated with longer-context evaluation is sufficient to produce
algorithmic length generalisation. ALiBi was introduced in a causal
language-modelling setting specifically to support evaluation on sequences
longer than those used during training~\cite{press2022train}, while RoPE
encodes relative positional relationships through rotations of query and key
representations~\cite{su2024roformer}. Previous work has also shown that NoPE
can exhibit competitive length-generalisation behaviour on some tasks
~\cite{kazemnejad2023impact}.

In the present experiments, RoPE provided a stable training baseline for
Delayed Copy but still collapsed at the first unseen length. It did not
provide an eligible baseline for Addition or Reverse. NoPE and the
bidirectional ALiBi adaptation did not learn any of the three training tasks
under the evaluated architecture and optimisation settings. These outcomes
show that the positional mechanism alone did not determine whether the model
would first learn the task or subsequently extrapolate.

The results should not be interpreted as a general contradiction of the
original ALiBi, RoPE, or NoPE studies. The present model was a compact
bidirectional encoder trained from scratch on deterministic
sequence-to-sequence tasks. In particular, its ALiBi implementation used a
symmetric absolute-distance bias rather than the causal setting evaluated in
the original work. The tasks, architecture, optimisation budget, and
attention structure therefore differ from those previous evaluations.

Instead, the comparison demonstrates an interaction between task structure
and positional treatment. Learned embeddings were eligible for Addition and
Reverse but seed unstable for Delayed Copy. Sinusoidal encoding was eligible
for Addition and Delayed Copy but seed unstable for Reverse. RoPE was eligible
only for Delayed Copy. No positional condition was consistently eligible
across all three tasks, and none of the eligible conditions extrapolated at
the sequence level.

The benchmark therefore supports the narrower conclusion that positional
encoding is one component of length-generalisation behaviour, not a sufficient
guarantee of it. A fuller comparison with prior work and the implications for
model design are developed in the Discussion chapter.

\section{Summary of Synthetic Findings}
\label{sec:synthetic-findings-summary}

The complete synthetic baseline matrix produced five eligible
task--encoding configurations, two seed-unstable configurations, and eight
configurations that did not satisfy the training-fit criterion. This
separation was necessary because poor performance at an unseen length has a
different interpretation depending on whether the model first learned the
training task.

All five eligible baselines exhibited immediate extrapolation collapse.
Addition with learned and sinusoidal encodings fell from mean training-length
exact-match accuracies of 92.0\% and 93.7\% to 0\% at length 32. Delayed
Copy with sinusoidal encoding and RoPE fell from 99.93\% and 100\% to 0\%
at length 256. Reverse with learned embeddings fell from 100\% to 0\% at
length 256. These outcomes were reproduced across all three seeds.

Token errors established that the sequence-level collapse was severe.
Addition token error increased to approximately 87\% at the first unseen
length and approached its 90\% uniform digit-prediction reference at longer
lengths. Delayed Copy and Reverse token errors increased to approximately
99\%, close to their uniform 100-symbol references. The 0\% exact-match
results therefore did not generally represent near-correct sequences
containing only occasional token errors.

The positional comparison did not identify one encoding that was consistently
successful across tasks. It instead revealed task-dependent training fit and
two cases of substantial seed instability. These findings reinforce the need
to establish a valid multi-seed baseline before interpreting extrapolation
performance.

This chapter has established whether failure occurred, its severity in
percentage terms, and the first sampled length at which it was observed. The
next chapter examines how the errors were distributed across output positions
and task-specific conditions, how attention statistics changed with length,
and whether the tested training interventions produced any observed
improvement.

% -----------------------------------------------------------------------------

\chapter{Failure Behaviour, Mitigation Strategies, and Taxonomy}
\label{chap:failure-analysis}

\noindent
The previous chapter established which configurations reliably learned the
training task and showed that all five eligible task--encoding baselines
collapsed at the first evaluated unseen length. This chapter moves beyond
the binary question of whether a model fails and examines how that failure
is expressed across the output sequence.

The analysis first examines position-wise token errors and task-specific
failure conditions. For Addition, errors are compared across output positions
and between digits with and without an incoming carry. For Delayed Copy,
errors are examined across the copied output sequence. For Reverse, errors
are analysed by output position and dependency distance. Attention statistics
are then used as descriptive evidence about changes in attention behaviour
across sequence lengths. Finally, the chapter compares mitigation variants
with their corresponding original baselines and organises the observed
behaviours into a failure taxonomy.

Because the detailed diagnostics use representative seed-42 checkpoints,
their purpose is to characterise the form of failure rather than to establish
multi-seed prevalence. Multi-seed baseline eligibility and robustness were
reported separately in Chapter~\ref{chap:synthetic-results}.

\section{Motivation for Failure Analysis}

Exact-match accuracy provides a strict sequence-level measure: an example is
counted as correct only when every evaluated output token is predicted
correctly. This is appropriate for algorithmic tasks, but once exact-match
accuracy reaches zero it does not reveal whether errors are concentrated in
particular output regions or distributed across the entire sequence.
Aggregate token accuracy provides additional information, but it can still
hide systematic variation across positions and task conditions.

The diagnostic analysis therefore uses four complementary sources of
evidence. First, position-wise token-error profiles identify whether failure
is concentrated near the beginning, middle, or end of the output. Second,
task-specific comparisons test whether Addition errors are associated with
incoming carries and whether Reverse errors vary with dependency distance.
Third, attention statistics describe how the learned attention distributions
change as the evaluated sequence length increases. Fourth, matched
baseline--mitigation comparisons test whether particular interventions
improve performance at unseen lengths while preserving training-length
performance.

These sources answer different questions and should not be treated as
interchangeable. Position-wise and task-specific diagnostics describe where
prediction errors occur. Attention statistics describe associated changes in
attention weights, but do not by themselves establish a causal or complete
mechanistic explanation of model behaviour~\cite{jain2019attention}.
Mitigation comparisons provide intervention evidence, although conclusions
remain limited to the particular configurations evaluated.

\section{Attention Diagnostic Metrics and Scope}

Attention diagnostics were computed from saved attention maps for selected
seed-42 checkpoints at the training length, the first evaluated unseen
length, and the longest test length. As defined in the methodology chapter,
the analysis uses normalised attention entropy, average attention distance,
and local attention ratio.

Normalised attention entropy indicates whether attention becomes more
concentrated or diffuse relative to the maximum entropy possible at a given
sequence length. Average attention distance summarises the distance between
query positions and the positions receiving attention. Local attention ratio
measures the proportion of attention assigned within a fixed neighbourhood
of each query position.

The statistics were averaged over attention heads, layers, and the examples
included in the diagnostic batches. The analysis used the first five
evaluation batches, corresponding to at most 40 Addition examples and five
Delayed Copy or Reverse examples. Consequently, these measurements are
treated as descriptive case studies rather than population-level estimates
or direct explanations of the prediction failures.

The position-wise and task-specific error analyses are presented before the
attention results because they provide the most direct evidence about the
observable form of failure.

\section{Position-Wise Failure Behaviour}
\label{sec:position-wise-failure}

The aggregate results establish that exact-match accuracy collapses at the
first evaluated unseen length, but they do not reveal whether token errors
are concentrated in particular parts of the output. To examine this, valid
target positions were divided into ten bins according to their relative
position in the output sequence. Token-error rates were then calculated
within each bin.

The analysis uses representative seed-42 checkpoints from the five
task--encoding configurations that satisfied the multi-seed baseline
eligibility criterion. The dotted reference lines indicate the expected
token-error rate under uniform prediction over the valid target vocabulary:
90\% for the ten Addition digits and 99\% for the 100-token Delayed Copy and
Reverse vocabularies.

\subsection{Addition}

\begin{figure}[htbp]
\centering
\includegraphics[width=\textwidth]
{addition_train_vs_first_ood_position_error.pdf}
\caption{Position-wise token-error rates for the eligible Addition
configurations at the training length and first evaluated unseen length.
Target positions are grouped into ten relative-position bins. The dotted
line denotes the 90\% uniform digit-prediction error reference. Results use
representative seed-42 checkpoints.}
\label{fig:addition-position-error}
\end{figure}

Figure~\ref{fig:addition-position-error} shows that both Addition models have
low token-error rates across the output at the training length. At length 32,
errors become distributed throughout the output, but they are not completely
uniform across position.

For the learned model, the overall token-error rate at length 32 was 88.0\%.
The error rate increased from 82.7\% in the start region to 86.5\% in the
middle and 89.9\% near the end, giving a positional range of 8.5 percentage
points. The sinusoidal model showed a stronger gradient: its error rate
increased from 77.1\% at the start to 83.9\% in the middle and 89.5\% at the
end, a range of 12.7 percentage points.

The failure is therefore broadly distributed rather than confined to a
single output region, but the lower error near the beginning indicates that
some limited positional structure remains at the first unseen length.
Performance approaches the uniform-prediction reference towards later output
positions.

\subsection{Delayed Copy}

\begin{figure}[htbp]
\centering
\includegraphics[width=\textwidth]
{copy_train_vs_first_ood_position_error.pdf}
\caption{Position-wise token-error rates for the eligible Delayed Copy
configurations at the training length and first evaluated unseen length.
Target positions are grouped into ten relative-position bins. The dotted
line denotes the 99\% uniform valid-token error reference. Results use
representative seed-42 checkpoints.}
\label{fig:copy-position-error}
\end{figure}

Figure~\ref{fig:copy-position-error} shows a different failure pattern.
Sinusoidal and RoPE have approximately zero token error throughout the output
at the training length. At length 256, however, both models produce error
rates close to the 99\% uniform-prediction reference across every relative
position.

The overall token-error rates were 99.0\% for both configurations. For the
sinusoidal model, the start-, middle-, and end-region error rates were
98.94\%, 98.97\%, and 99.10\%, respectively. For RoPE, the corresponding
values were 99.00\%, 98.90\%, and 99.05\%. The positional ranges were only
approximately 0.23 and 0.22 percentage points.

Delayed Copy therefore exhibits an almost spatially uniform near-reference
collapse. The model does not retain a clearly better-performing region near
the beginning or end of the copied output when the sequence length doubles.

\subsection{Reverse}

\begin{figure}[htbp]
\centering
\includegraphics[width=\textwidth]
{reverse_train_vs_first_ood_position_error.pdf}
\caption{Position-wise token-error rates for Reverse with learned positional
embeddings at the training length and first evaluated unseen length. Target
positions are grouped into ten relative-position bins. The dotted line
denotes the 99\% uniform valid-token error reference. Results use the
representative seed-42 checkpoint.}
\label{fig:reverse-position-error}
\end{figure}

As shown in Figure~\ref{fig:reverse-position-error}, the learned Reverse model
has zero token error across the output at the training length. At length 256,
the overall token-error rate rises to 98.79\%, and errors occur throughout
the output sequence.

The start-, middle-, and end-region error rates were 98.87\%, 98.85\%, and
98.30\%, respectively. Their range was only 0.71 percentage points.
Consequently, Reverse failure is not concentrated at a particular relative
output position. Instead, the model changes from uniformly correct behaviour
at the training length to an almost uniformly incorrect output at the first
unseen length.

\subsection{Comparison Across Tasks}

The position-wise diagnostics distinguish two forms of spatial behaviour.
Addition retains a modest positional gradient at the first unseen length,
with lower errors near the beginning and errors approaching the
uniform-prediction reference towards the end. Delayed Copy and Reverse show
much more uniform failure, with error rates close to their 99\% reference
throughout the output.

These results demonstrate why exact-match accuracy alone is insufficient.
Although all five eligible configurations have 0\% exact-match accuracy at
the first unseen length, the distribution of their token errors is not
identical.

\section{Task-Specific Failure Conditions}
\label{sec:task-specific-failure}

Position-wise profiles identify where errors occur in the output, but they do
not test whether those errors correspond to the computation required by each
task. Two additional analyses were therefore performed. Addition errors were
conditioned on the presence of an incoming carry, while Reverse errors were
conditioned on the distance between each output position and its corresponding
source position.

\subsection{Addition: Incoming-Carry Condition}

\begin{figure}[htbp]
\centering
\includegraphics[width=\textwidth]
{addition_carry_condition_error.pdf}
\caption{Addition error rates for ordinary target digits with and without an
incoming carry. The learned and sinusoidal seed-42 checkpoints are shown
across all evaluated lengths. The optional leading result digit is excluded
from this comparison. The dotted line denotes the 90\% uniform
digit-prediction error reference.}
\label{fig:addition-carry-error}
\end{figure}

Figure~\ref{fig:addition-carry-error} compares ordinary target digits according
to whether their calculation receives a carry from the preceding digit
column. At the first unseen length, the learned model produced an error rate
of 88.14\% without an incoming carry and 89.09\% with an incoming carry. The
difference was 0.96 percentage points. For the sinusoidal model, the
corresponding error rates were 88.46\% and 88.46\%, differing by less than
0.01 percentage points.

The two conditions remain closely aligned at the longer evaluation lengths.
Therefore, the observed Addition collapse is not primarily concentrated on
digits requiring an incoming carry. This does not show that the models
learned a length-general carry procedure; rather, it shows that once
extrapolation fails, ordinary digits with and without incoming carries are
both predicted poorly.

\subsection{Reverse: Dependency Distance}

For an output position $i$ in a sequence of length $L$, the corresponding
source position in the Reverse task is $L-1-i$. The normalised dependency
distance used in this analysis is therefore

\[
d_i =
\frac{\left|(L-1-i)-i\right|}{L-1}.
\]

A value near zero represents a relatively short input--output dependency,
whereas a value near one represents a dependency spanning almost the full
sequence.

\begin{figure}[htbp]
\centering
\includegraphics[width=0.88\textwidth]
{reverse_dependency_distance_error.pdf}
\caption{Reverse token-error rates grouped by normalised input--output
dependency distance for the representative learned seed-42 checkpoint.
Distances are divided into ten equal-width bins.}
\label{fig:reverse-dependency-error}
\end{figure}

Figure~\ref{fig:reverse-dependency-error} shows zero error across all
dependency-distance bins at the training length. At the first unseen length,
error rates range only from approximately 98.54\% to 98.92\% across the ten
bins. Similarly high and nearly uniform error rates occur at lengths 512 and
1024.

The failure is therefore not concentrated only on output tokens requiring the
longest dependencies. Shorter and longer reverse mappings fail at comparable
rates once the model is evaluated beyond its training length. This supports
the interpretation of a distributed reversal failure rather than a gradual
loss of accuracy as dependency distance increases.

\section{Attention Behaviour Across Tasks}
\label{sec:attention-behaviour}

Attention statistics provide descriptive evidence about changes in model
behaviour across sequence lengths. Table~\ref{tab:attention-summary} reports
normalised attention entropy and local attention ratio for the representative
seed-42 checkpoints. Values are descriptive means across the two Transformer
layers and the first five diagnostic batches; they are not multi-seed
estimates.

\begin{table}[htbp]
\centering
\small
\setlength{\tabcolsep}{4pt}
\begin{tabular}{llcccccc}
\hline
& &
\multicolumn{3}{c}{Normalised entropy} &
\multicolumn{3}{c}{Local attention ratio} \\
\cline{3-5}\cline{6-8}
Task & Encoding &
Train & First OOD & Longest &
Train & First OOD & Longest \\
\hline
Addition & Learned
& 0.157 & 0.278 & 0.594
& 0.496 & 0.280 & 0.012 \\
Addition & Sinusoidal
& 0.345 & 0.475 & 0.588
& 0.534 & 0.360 & 0.012 \\
Delayed Copy & Sinusoidal
& 0.645 & 0.627 & 0.630
& 0.075 & 0.032 & 0.007 \\
Delayed Copy & RoPE
& 0.605 & 0.577 & 0.638
& 0.069 & 0.027 & 0.005 \\
Reverse & Learned
& 0.539 & 0.731 & 0.910
& 0.119 & 0.073 & 0.022 \\
\hline
\end{tabular}
\caption{Attention statistics for representative seed-42 checkpoints at the
training length, first evaluated out-of-distribution (OOD) length, and longest
test length. The fixed local neighbourhood is defined by
$|i-j|\leq 10$. Values are descriptive means over two layers and five
diagnostic batches.}
\label{tab:attention-summary}
\end{table}

\subsection{Addition}

For Addition, normalised entropy increased with sequence length for both
eligible encodings. For the learned model, it increased from 0.157 at length
16 to 0.278 at length 32 and 0.594 at length 1024. The sinusoidal model showed
a similar increase from 0.345 to 0.475 and then 0.588. At the same time, the
local attention ratio decreased from approximately 0.50 at the training
length to approximately 0.01 at the longest length.

The first entropy increase coincided with exact-match accuracy falling from
94.1\% to 0\% for the learned seed-42 model and from 94.7\% to 0\% for the
sinusoidal model. This establishes an association between extrapolation
failure and a redistribution of attention in these checkpoints. It does not,
however, establish that increasing entropy caused the errors. The
position-wise and carry-condition results provide more direct evidence about
how Addition predictions failed.

\subsection{Delayed Copy}

Delayed Copy produced a different pattern. Normalised entropy remained
approximately stable for both eligible encodings. For the sinusoidal model,
the values were 0.645, 0.627, and 0.630 at lengths 128, 256, and 1024,
respectively. For RoPE, they were 0.605, 0.577, and 0.638. Exact-match
accuracy nevertheless fell from 99.8\% and 100\% at the training length to
0\% at length 256.

The absence of a substantial entropy increase at the first unseen length is
important. It shows that immediate extrapolation collapse was not always
accompanied by greater attention diffusion. The local attention ratio
decreased for both encodings, but this statistic uses a fixed
$\pm 10$-position neighbourhood. Its value is therefore expected to decrease
partly as the total sequence length increases and must be interpreted
cautiously.

\subsection{Reverse}

For Reverse with learned positional embeddings, normalised entropy increased
from 0.539 at length 128 to 0.731 at length 256 and 0.910 at length 1024.
The local attention ratio simultaneously decreased from 0.119 to 0.073 and
then 0.022. The first change coincided with exact-match accuracy falling from
100\% to 0\%.

These results show that the representative Reverse model's attention became
more diffuse outside the training length. However, the position-wise and
dependency-distance analyses show whether the resulting prediction errors
were spatially concentrated. The attention measurements alone cannot
determine whether the model learned a length-specific reversal mechanism.

\subsection{Interpretation Across Tasks}

Average attention distance increased with test length in every task.
However, this metric is measured in raw token positions, so part of the
increase follows mechanically from the larger set of possible distances in
longer sequences. It is therefore not treated as independent evidence of
improved long-range reasoning or of failure.

Taken together, the attention diagnostics do not support a single universal
attention-based explanation. Addition and Reverse showed increased
normalised entropy, whereas Delayed Copy collapsed without a comparable
increase at the first unseen length. Attention change is therefore retained
as descriptive supporting evidence rather than used as a separate failure
category in the final taxonomy.

\section{Mitigation Strategy Design}
\label{sec:mitigation-design}

The mitigation experiments test specific hypotheses about why a model that
learns the training task might fail at longer lengths. The purpose is not
simply to train additional variants, but to determine whether changing the
distribution or ordering of training lengths alters the observed
extrapolation behaviour.

The experiments use representative seed-42 configurations. A mitigation is
considered interpretable as a length-generalisation test only when it first
preserves training-length performance. The same 90\% exact-match eligibility
threshold used for the baseline analysis is therefore applied to the
mitigation variants. If a variant falls below this threshold, its
out-of-distribution results are classified as inconclusive because the
training task was not learned sufficiently.

\begin{table}[htbp]
\centering
\small
\begin{tabular}{p{2.7cm}p{3.2cm}p{7.0cm}}
\hline
Condition & Applied setting & Specific question tested \\
\hline
Single-length control
& Delayed Copy, sinusoidal
& Does the mitigation training pipeline reproduce the original result when
the length distribution remains unchanged? \\

Mixed-length training
& Addition, learned; Delayed Copy, sinusoidal
& Does simultaneous exposure to several lengths reduce specialisation to one
fixed training length? Addition uses lengths 4, 8, 12, and 16; Delayed Copy
uses lengths 32, 64, 96, and 128. \\

Strengthened mixed-length training
& Addition, learned
& Was the initial mixed-length result caused by insufficient data or
optimisation time? The second variant increases both the number of training
examples and the number of epochs. \\

Curriculum training
& Delayed Copy, sinusoidal and RoPE; Reverse, learned
& Does sequential training at lengths 32, 64, 96, and 128 reduce an
optimisation barrier associated with training directly at length 128? \\

Randomised padded training
& Delayed Copy, sinusoidal and RoPE
& Does combining variable-length examples in a shuffled, right-padded dataset
reduce reliance on fixed separator and output positions? This intervention
varies effective sequence length but does not randomly shift tokens to new
absolute offsets. \\
\hline
\end{tabular}
\caption{Mitigation conditions and the specific hypotheses tested by each
intervention.}
\label{tab:mitigation-hypotheses}
\end{table}

The primary comparison is exact-match accuracy at the training length and
the first evaluated unseen length. The longest test length is also reported,
but it is not used alone to judge mitigation success. A useful mitigation
would preserve training-length performance and produce a meaningful
improvement over the corresponding baseline at unseen lengths.

\section{Mitigation Results}
\label{sec:mitigation-results}

The mitigation experiments produced two principal outcomes. Some variants
failed to learn the training task and therefore did not provide a valid test
of extrapolation. Other variants preserved training-length performance but
showed no observed improvement at the first unseen or longest evaluated
length.

These experiments use one representative seed and should consequently be
interpreted as targeted intervention tests rather than multi-seed estimates.

\subsection{Addition: Mixed-Length Training}

\begin{figure}[htbp]
\centering
\includegraphics[width=0.95\textwidth]
{addition_learned_baseline_vs_mitigation.pdf}
\caption{Exact-match accuracy for the original Addition learned baseline and
two mixed-length variants. Results use seed 42. Both mixed-length variants
fall below the training-length eligibility threshold, so their extrapolation
results are classified as inconclusive.}
\label{fig:addition-mitigation}
\end{figure}

Figure~\ref{fig:addition-mitigation} compares the original learned-position
baseline with two mixed-length variants. The original baseline achieved
94.1\% exact-match accuracy at length 16 and 0\% at the first unseen length
of 32.

The first mixed-length variant achieved 0\% exact-match accuracy at length
16. The strengthened variant, which used more examples and additional
epochs, also achieved 0\% at the training length. Because neither variant
learned the length-16 task, their 0\% performance at longer lengths cannot be
interpreted as evidence about extrapolation.

The appropriate conclusion is therefore not that mixed-length training was
ineffective at improving an otherwise learned model. Instead, mixed-length
training introduced a training-length optimisation failure under both
evaluated settings. Whether a differently tuned mixed-length procedure could
preserve the baseline fit remains unresolved.

\subsection{Delayed Copy with Sinusoidal Encoding}

\begin{figure}[htbp]
\centering
\includegraphics[width=0.95\textwidth]
{copy_sinusoidal_baseline_vs_mitigation.pdf}
\caption{Exact-match accuracy for the original Delayed Copy sinusoidal
baseline, the single-length control, and the mixed-length, curriculum, and
randomised-padded variants. Results use seed 42; no multi-seed error bars are
shown.}
\label{fig:copy-sinusoidal-mitigation}
\end{figure}

The single-length control reproduced the original sinusoidal baseline exactly,
with a maximum difference of 0 percentage points across the evaluated
lengths. Both achieved 99.8\% exact-match accuracy at length 128 and 0\% at
lengths 256, 512, and 1024. This confirms that the mitigation pipeline could
reproduce the reference failure curve when the training-length distribution
was unchanged.

The mixed-length variant achieved 0\% exact-match accuracy at the training
length and is therefore classified as training-length underfitting. In
contrast, curriculum and randomised-padded training both achieved 100\%
exact-match accuracy at length 128. These two variants provide valid
extrapolation tests because they preserved training-length performance.

Neither valid mitigation improved performance at the first unseen length.
Curriculum and randomised-padded training both achieved 0\% exact-match
accuracy at lengths 256, 512, and 1024, matching the extrapolation behaviour
of the original baseline. Thus, gradual length exposure and jointly padded
variable-length exposure produced no observed exact-match improvement in
this configuration.

\subsection{Delayed Copy with RoPE}

\begin{figure}[htbp]
\centering
\includegraphics[width=0.95\textwidth]
{copy_rope_baseline_vs_mitigation.pdf}
\caption{Exact-match accuracy for the original Delayed Copy RoPE baseline and
the curriculum and randomised-padded variants. All conditions achieve 100\%
at the training length but 0\% at every evaluated unseen length. Results use
seed 42; no multi-seed error bars are shown.}
\label{fig:copy-rope-mitigation}
\end{figure}

Figure~\ref{fig:copy-rope-mitigation} shows that the original baseline,
curriculum variant, and randomised-padded variant all achieved 100\%
exact-match accuracy at length 128. Training-length underfitting therefore
does not explain the mitigation outcomes for this configuration.

All three conditions nevertheless achieved 0\% exact-match accuracy at
lengths 256, 512, and 1024. The two interventions preserved the ability to
solve the training task but produced no observed improvement at unseen
lengths. This result is consistent with the sinusoidal experiment and shows
that the outcome was not limited to one positional encoding.

\subsection{Reverse Curriculum Qualification}

The Reverse curriculum run achieved 100\% exact-match accuracy at length 128
and 0\% at the evaluated unseen lengths. However, this comparison is not
fully matched: the original Reverse baseline used a vocabulary size of 100,
whereas the curriculum configuration used a vocabulary size of 102.
Consequently, the run is treated as exploratory supporting evidence rather
than as a controlled baseline--mitigation comparison. A corrected rerun
would be required before making a stronger claim about curriculum training
on Reverse.

\subsection{Interpretation of the Mitigation Evidence}

The mitigation results do not support the broad claim that exposure to varied
training lengths necessarily improves extrapolation. Mixed-length variants
underfit in the evaluated Addition and sinusoidal Delayed Copy settings.
Curriculum and randomised-padded Delayed Copy variants retained
training-length performance but showed no exact-match improvement at the
first unseen length.

These are bounded negative results. They show that the specific
interventions, configurations, and representative seed evaluated here did
not improve exact-match extrapolation. They do not establish that the
failures are universally resistant to mitigation or that more extensive
hyperparameter tuning, alternative curricula, or different model
architectures could not improve performance.

\section{Proposed Failure Taxonomy}
\label{sec:failure-taxonomy}

The failure taxonomy is the main interpretive contribution of this project.
Its purpose is to prevent distinct experimental outcomes from being grouped
under the general statement that a model ``fails at longer lengths''. The
categories separate failure to learn the original task, instability across
random seeds, failure after a valid training fit, the spatial form of the
resulting errors, and the outcome of a tested intervention.

\begin{table}[htbp]
\centering
\small
\begin{tabular}{p{3.1cm}p{4.6cm}p{5.3cm}}
\hline
Failure category & Operational evidence & Interpretation \\
\hline
Training-length underfitting
& Training-length exact-match accuracy is below 90\%.
& The task was not learned sufficiently, so longer-length performance cannot
be used as evidence of extrapolation failure. \\

Seed instability
& Only one or two of the three seeds reach the 90\% training-fit threshold,
or training performance varies substantially across seeds.
& A pooled mean can conceal qualitatively different runs; successful and
unsuccessful seeds should be reported separately. \\

Immediate extrapolation collapse
& All three seeds satisfy baseline eligibility, but mean exact-match accuracy
falls below 10\% at the first evaluated unseen length.
& Training-length success does not transfer to the first sampled longer
length. ``Immediate'' refers to the evaluated length grid and does not locate
the precise failure point between evaluations. \\

Spatial error structure
& Position-wise or task-conditioned token-error rates vary across the target,
carry condition, or dependency-distance bins.
& This category describes how failure is distributed. Addition retains a
modest positional gradient, whereas Delayed Copy and Reverse show almost
uniform near-reference error. \\

No observed improvement under tested mitigation
& A mitigation preserves training-length accuracy but does not improve
exact-match accuracy at the first unseen length relative to the reference
baseline.
& The evaluated intervention did not improve extrapolation in the tested
configuration. This is a scoped experimental result, not evidence of
universal resistance to mitigation. \\
\hline
\end{tabular}
\caption{Failure taxonomy derived from baseline eligibility, multi-seed
robustness, position-wise diagnostics, and mitigation experiments.}
\label{tab:failure-taxonomy}
\end{table}

The taxonomy is applied sequentially. Training-length performance is checked
first because a model must learn the original task before its extrapolation
can be evaluated. Multi-seed consistency is then examined to determine
whether the baseline is reliable. For eligible configurations, the first
unseen length identifies whether immediate extrapolation collapse occurs.
Position-wise and task-specific analyses subsequently characterise the
spatial form of that collapse. Finally, mitigation outcomes are evaluated
only when the intervention preserves training-length performance.

This ordering prevents several misleading comparisons. For example, NoPE and
ALiBi frequently perform poorly in the present setup, but their outcomes are
classified as training-length underfitting rather than extrapolation
collapse. Delayed Copy with learned positional embeddings and Reverse with
sinusoidal encoding are classified as seed-unstable because only some seeds
learn the training task. In contrast, the five eligible task--encoding
configurations provide valid evidence of immediate extrapolation collapse.

Attention behaviour is not used as a separate failure category. Addition and
Reverse show increasing normalised attention entropy, but Delayed Copy
collapses without a comparable entropy increase at the first unseen length.
Attention statistics are therefore retained as descriptive supporting
evidence rather than treated as a universal failure mechanism.

\section{Summary}

This chapter examined how the eligible models fail after leaving the training
length. The position-wise analysis showed that Addition retains a modest
error gradient, with lower errors near the beginning and near-reference
errors towards later output positions. Delayed Copy and Reverse instead show
almost uniform near-reference error across the output.

The task-specific diagnostics further showed that Addition errors were not
substantially concentrated on digits with an incoming carry and that Reverse
errors did not increase meaningfully with dependency distance. Attention
statistics produced task-dependent patterns: entropy increased for Addition
and Reverse but remained approximately stable for Delayed Copy at the first
unseen length. Attention diffusion therefore does not provide a universal
explanation of the observed failures.

The mitigation experiments also produced distinct outcomes. Mixed-length
variants underfit in the evaluated settings, while curriculum and
randomised-padded Delayed Copy variants preserved training-length performance
but produced no observed improvement at unseen lengths. These results were
organised into five categories: training-length underfitting, seed
instability, immediate extrapolation collapse, spatial error structure, and
no observed improvement under tested mitigation.

The taxonomy should be interpreted as an evidence-based framework for the
experiments in this dissertation rather than as an exhaustive taxonomy of
all possible length-generalisation failures. It provides the basis for the
broader interpretation, limitations, and recommendations developed in the
following chapters.
% ----------------------------------------------------------------------

\chapter{Exploratory Real-World Extension}
\label{chap:real-world-extension}

\noindent
The main experiments use synthetic algorithmic tasks because they provide
controlled tests of length extrapolation and permit exact analysis of output
errors. However, these tasks are deliberately simplified and do not represent
all properties of applied sequential data. This chapter presents a limited
real-world extension using the SelfRegulationSCP2 multivariate time-series
classification dataset.

The purpose of this experiment is to provide an exploratory contrast rather
than to validate the synthetic findings or establish state-of-the-art
performance. In particular, classification accuracy and sequence-level
exact-match accuracy measure different forms of performance and should not be
compared directly.

\section{Dataset and Task}

SelfRegulationSCP2 is part of the UEA multivariate time-series classification
archive~\cite{bagnall2018uea}. It contains EEG recordings from a binary
classification task involving the self-regulation of slow cortical
potentials. The predefined split contains 200 training examples and 180 test
examples. Each example contains seven channels and 1152 time points.

The model predicts one class label for each input sequence. This differs from
the synthetic tasks, where a prediction is counted as an exact match only when
every valid output token is correct. A sequence-level classifier may use
partial temporal evidence, while the synthetic tasks require a complete
structured output.

\begin{table}[htbp]
\centering
\begin{tabular}{ll}
\hline
Component & Setting \\
\hline
Dataset & SelfRegulationSCP2 \\
Data type & Multivariate EEG time series \\
Training examples & 200 \\
Test examples & 180 \\
Input channels & 7 \\
Classes & 2 \\
Full sequence length & 1152 \\
Training prefix length & 256 \\
Evaluation prefix lengths & 256, 512, 1024, 1152 \\
Metric & Classification accuracy \\
Positional encodings & Learned, sinusoidal, NoPE, ALiBi, RoPE \\
Seeds & 42, 123, 2024 \\
\hline
\end{tabular}
\caption{Dataset and experimental configuration for the exploratory
SelfRegulationSCP2 extension.}
\label{tab:real-world-setup}
\end{table}

\section{Experimental Protocol}

The experiment follows a train-short/test-long prefix protocol. Each
Transformer classifier is trained using the first 256 time points of every
training example. The resulting checkpoint is evaluated on test prefixes of
length 256, 512, 1024, and 1152. The test examples remain the same across
lengths; only the amount of each time series supplied to the classifier
changes.

Continuous input channels are projected from seven dimensions to the
128-dimensional Transformer representation using a linear layer. The model
then applies the same two-block Transformer structure used in the synthetic
experiments. Mean pooling across the temporal dimension produces a
sequence-level representation, which is passed to a two-class linear
classification head.

Normalisation statistics are calculated using only the training split and
then applied to both training and test data. Each configuration is trained
for 50 epochs using a batch size of 16 and repeated across three random
seeds. Results are reported as mean test accuracy, with error bars denoting
one standard error across the three seeds.

\section{Real-World Results}

The existing real-world plot should be regenerated in the final visual style
before compilation. The revised figure will be saved as
\texttt{selfregulationscp2\_accuracy\_mean\_se.pdf}.

\begin{figure}[htbp]
\centering
\includegraphics[width=0.92\textwidth]
{selfregulationscp2_accuracy_mean_se.pdf}
\caption{Mean SelfRegulationSCP2 test accuracy across three seeds as the
evaluation prefix length increases. Error bars show $\pm 1$ standard error.
The dotted horizontal line denotes the 50\% uniform two-class prediction
reference, and the vertical marker identifies the training prefix length of
256. Markers are offset slightly horizontally to reveal coincident estimates;
axis ticks show the nominal prefix lengths.}
\label{fig:selfregulationscp2-accuracy}
\end{figure}

\begin{table}[htbp]
\centering
\small
\begin{tabular}{lccc}
\hline
Encoding &
Accuracy at 256 (\%) &
Accuracy at 1152 (\%) &
Change (percentage points) \\
\hline
Learned     & $49.44 \pm 0.64$ & $53.52 \pm 0.37$ & $+4.07$ \\
Sinusoidal  & $52.04 \pm 2.73$ & $53.52 \pm 3.43$ & $+1.48$ \\
NoPE        & $52.96 \pm 1.13$ & $50.37 \pm 3.96$ & $-2.59$ \\
ALiBi       & $50.19 \pm 0.98$ & $51.48 \pm 0.98$ & $+1.30$ \\
RoPE        & $50.93 \pm 0.74$ & $49.07 \pm 0.37$ & $-1.85$ \\
\hline
\end{tabular}
\caption{Mean classification accuracy and standard error at the training
prefix length and full sequence length. Changes are calculated from the
unrounded mean accuracies.}
\label{tab:real-world-results}
\end{table}

Figure~\ref{fig:selfregulationscp2-accuracy} shows that mean accuracy remains
within approximately 49--56\% across the evaluated lengths. No positional
encoding exhibits the immediate transition from high training-length
performance to zero accuracy observed in the eligible synthetic baselines.

At length 512, RoPE produces the highest individual mean of 55.56\%, but its
mean accuracy decreases to 49.07\% at length 1152. At the full length, learned
and sinusoidal encodings have the highest mean accuracy at 53.52\%. However,
these differences are small relative to the overall performance level and
seed variation, and no statistical significance test was conducted.
Consequently, the experiment does not establish a reliable ranking among the
positional encodings.

The most important qualification is that held-out accuracy at the training
prefix length is already modest. Mean accuracy at length 256 ranges only from
49.44\% to 52.96\%, close to the 50\% uniform two-class reference. Therefore,
the absence of a sharp decline cannot be interpreted as evidence that the
models learned a robust classifier that generalises across length. Unlike the
eligible synthetic baselines, the real-world experiment does not begin from
strong held-out task performance.

\section{Comparison with the Synthetic Benchmark}

The real-world curve differs descriptively from the synthetic exact-match
curves. Eligible synthetic configurations achieve at least 90\% exact-match
accuracy at the training length and then fall to 0\% at the first evaluated
unseen length. The SelfRegulationSCP2 models instead remain within a narrow,
modest accuracy range as additional time points are supplied.

Several factors prevent a causal interpretation of this contrast. First,
classification produces one label, whereas the synthetic tasks require an
entire output sequence. Second, the evaluation metrics have substantially
different sensitivities to partial correctness. Third, the
SelfRegulationSCP2 models do not establish a strong held-out reference
performance at length 256. Finally, longer prefixes introduce additional
signal that may be informative, irrelevant, or noisy depending on the
example.

The experiment therefore shows that catastrophic exact-match collapse is not
visible in this particular classification curve. It does not show that the
classifier learned a length-general computation, nor does it establish that
real-world classification tasks are generally less affected by sequence
length.

\section{Limitations of the Extension}

This extension has several important limitations. It uses only one relatively
small dataset, with 200 training examples and 180 test examples. No specialised
time-series baseline or state-of-the-art architecture is included, and the
Transformer hyperparameters were selected primarily for comparability with
the synthetic experiments rather than optimised for this dataset.

The models are evaluated only through classification accuracy. No
position-wise output diagnostic is possible because the task produces one
label per sequence. In addition, the three-seed results provide only a small
estimate of variability, and no formal significance testing or confidence
interval analysis was performed.

Most importantly, test accuracy at the training prefix length remains close
to the uniform two-class reference. The experiment therefore cannot cleanly
separate failure to generalise across examples from failure to generalise
across input length. A stronger extension would first establish a competitive
held-out baseline and would then examine whether that performance is
preserved when the input-length distribution changes.

Future work should include additional long-sequence datasets, majority-class
and specialised time-series baselines, task-specific hyperparameter
selection, and a larger number of seeds.

\section{Summary}

The SelfRegulationSCP2 extension provides an exploratory contrast to the
controlled synthetic benchmark. Mean classification accuracy remains within
approximately 49--56\% as the input prefix increases from 256 to 1152, rather
than showing an immediate fall to zero.

However, performance is already close to the uniform two-class reference at
the training prefix length. The experiment therefore provides evidence of a
flat but weak classification curve, not successful real-world length
generalisation. Its main value is to demonstrate why conclusions about length
behaviour depend on baseline validity, task structure, output type, and
evaluation metric.

%=========================================================================

\chapter{Discussion}
\label{chap:discussion}

\noindent
This chapter interprets the experimental findings in relation to the primary
research question and its five supporting questions. The central issue is not
only whether performance decreases at longer lengths, but whether the
observed outcome represents failure to learn the training task, instability
across seeds, collapse after a valid training fit, a particular spatial error
pattern, or no observed improvement under a tested intervention.

The results demonstrate that these distinctions materially change how poor
long-length accuracy should be interpreted. They also show that neither a
particular positional encoding nor exposure to varied training lengths is,
by itself, sufficient to guarantee length-general behaviour in the evaluated
configurations.

\section{Answers to the Research Questions}
\label{sec:discussion-research-questions}

\subsection{RQ1: Which Configurations Reliably Learn the Training Task?}

Five task--encoding configurations satisfied the baseline-eligibility
criterion across all three seeds: Addition with learned and sinusoidal
positional encodings, Delayed Copy with sinusoidal encoding and RoPE, and
Reverse with learned positional embeddings. These configurations provide
valid starting points for analysing extrapolation because they reliably
learned the corresponding training-length task.

Other configurations produced qualitatively different outcomes. Delayed Copy
with learned positional embeddings and Reverse with sinusoidal encoding
succeeded for only some seeds. Pooling these successful and unsuccessful runs
into a single mean would conceal the fact that the configurations were
unstable at the training length. NoPE and ALiBi, together with several other
task--encoding combinations, did not reliably learn the training task under
the present architecture and optimisation settings.

This result supports the decision to treat baseline eligibility as a
prerequisite for extrapolation analysis. More generally, reporting multiple
runs is important for stochastic neural models because conclusions based on
one seed may not represent the distribution of possible outcomes
\cite{reimers2017reporting}. In this project, seed variation is not merely
uncertainty around a common mean; in some configurations it separates
successful task learning from complete training-length failure.

\subsection{RQ2: When Does Eligible Performance Collapse?}

All five eligible configurations fell to 0\% exact-match accuracy at the
first evaluated unseen length. Addition collapsed when the digit length
increased from 16 to 32, while Delayed Copy and Reverse collapsed when the
sequence length increased from 128 to 256. Performance remained at 0\%
exact match at all longer evaluated lengths.

The behaviour is therefore classified as immediate extrapolation collapse.
However, ``immediate'' is defined relative to the sampled evaluation grid.
The experiments establish that collapse occurred by lengths 32 and 256,
respectively, but do not identify the precise transition point between the
training length and the first evaluated unseen length.

This result is consistent with previous evidence that strong performance on
short reasoning problems does not necessarily transfer to longer instances
\cite{anil2022exploring}. It is also compatible with the argument that
Transformer length generalisation depends on whether optimisation discovers
a computation whose structure remains valid as length changes
\cite{zhou2024algorithms}.

The results do not prove that a particular internal algorithm was learned at
the training length. They establish the more limited empirical conclusion
that the behaviour producing high training-length accuracy did not preserve
complete sequence-level correctness at the first sampled longer length.

\subsection{RQ3: Where Within the Output Do Errors Occur?}

The position-wise diagnostics show that 0\% exact-match accuracy does not
correspond to an identical token-error pattern across tasks. At the first
unseen length, Addition showed a modest positional gradient. Errors were
lower near the beginning of the output and increased towards later positions,
where they approached the 90\% uniform digit-prediction reference. The
sinusoidal model showed a larger positional range than the learned model.

Delayed Copy and Reverse exhibited a different form of failure. Their
token-error rates were close to the 99\% uniform valid-token reference across
almost every relative output position. These configurations changed from
approximately uniform correctness at the training length to approximately
uniform incorrectness at the first unseen length.

The task-specific diagnostics further narrow the interpretation. For
Addition, ordinary digits with and without an incoming carry had almost
identical error rates. At the first unseen length, their difference was less
than one percentage point for the learned model and effectively zero for the
sinusoidal model. The failure was therefore not concentrated specifically on
carry-dependent digits.

For Reverse, error rates were similarly high across short and long
normalised input--output dependency distances. The failure cannot therefore
be described as a progressive inability to handle only the most distant
reverse mappings. Once the model left the training length, mappings across
the full dependency-distance range failed at comparable rates.

These findings are important because they replace a general statement that
the models ``fail'' with a description of failure behaviour. Addition retains
some position-dependent structure, whereas Delayed Copy and Reverse show
distributed near-reference collapse.

\subsection{RQ4: How Does Positional Encoding Affect the Outcome?}

Positional encoding strongly affected whether the model learned the
training-length task, but no encoding was uniformly successful across all
three tasks. Learned positional embeddings were eligible for Addition and
Reverse but seed-unstable for Delayed Copy. Sinusoidal encoding was eligible
for Addition and Delayed Copy but seed-unstable for Reverse. RoPE was eligible
for Delayed Copy but underfit or only partially fit the other tasks.

Consequently, the experiments do not support a universal ranking of
positional encodings. Their effects depend on the task, architecture,
attention mask, and optimisation procedure. This task dependence is
consistent with broader evidence that positional encoding is an important
but context-dependent factor in Transformer length generalisation
\cite{kazemnejad2023impact}.

The underfitting of NoPE and ALiBi should not be interpreted as a general
failure of those methods. Previous studies have reported effective
length-related behaviour for NoPE in decoder-only Transformers
\cite{kazemnejad2023impact} and for ALiBi in causal language modelling
\cite{press2022train}. The present model is instead a compact,
encoder-style Transformer using bidirectional attention. Its ALiBi
implementation applies a symmetric distance bias rather than the original
causal formulation. These architectural differences prevent a direct
method-level comparison.

Similarly, RoPE provides a systematic rotation of query and key
representations \cite{su2021roformer}, but this representational mechanism
does not by itself guarantee that optimisation will discover a
length-general task solution. RoPE fitted Delayed Copy consistently and still
collapsed at the first unseen length.

\subsection{RQ5: Do the Tested Mitigations Improve Extrapolation?}

The mitigation experiments produced two distinct outcomes. The Addition
mixed-length variants and the sinusoidal Delayed Copy mixed-length variant
underfit at the training length. Their long-length results are therefore
inconclusive as tests of extrapolation.

Curriculum and randomised-padded training preserved training-length
performance for the eligible Delayed Copy configurations. Nevertheless, they
produced 0\% exact-match accuracy at the first unseen length, matching the
original baselines. These constitute valid negative results: the
interventions preserved the training task but produced no observed
extrapolation improvement in the tested settings.

The curriculum outcome does not show that curriculum learning is generally
ineffective. Curriculum learning is a broad training principle in which
examples are presented in an intended order of difficulty
\cite{bengio2009curriculum}. The present experiment tests only one
four-stage length curriculum and one representative seed.

The randomised-padded experiment must also be distinguished from randomised
positional encoding. The intervention used here combines shuffled,
variable-length examples in a right-padded dataset. It does not sample
ordered positional indices from a larger range. It therefore does not
implement the method proposed by Ruoss et al., who directly randomise
positional indices to expose the model to a broader positional distribution
\cite{ruoss2023randomized}. The negative result reported here does not
contradict their findings.

The evidence supports the scoped category ``no observed improvement under
tested mitigation''. It does not support the stronger claim that the
failures are inherently resistant to mitigation.

\section{Contribution of the Failure Taxonomy}
\label{sec:discussion-taxonomy}

The main interpretive contribution is a five-part taxonomy based on observable
experimental evidence:

\begin{enumerate}
    \item training-length underfitting;
    \item seed instability;
    \item immediate extrapolation collapse;
    \item spatial error structure; and
    \item no observed improvement under tested mitigation.
\end{enumerate}

These categories are not intended to represent mutually exclusive causal
mechanisms. Instead, they form a sequential framework for interpreting
experimental outcomes.

Training-length performance is examined first because extrapolation cannot be
meaningfully assessed before the original task is learned. Seed consistency
is then checked because a high pooled mean may hide failed runs. Eligible
configurations are evaluated across unseen lengths to identify the first
observed collapse. Position-wise and task-specific diagnostics characterise
how errors are distributed. Finally, an intervention is judged only after
checking whether it preserves the training baseline.

This structure provides more information than a single long-length accuracy
table. For example, an ALiBi configuration that never learns the training task
and a learned-position configuration that moves from 100\% to 0\% exact match
may have the same accuracy at the longest length, but they do not provide the
same evidence about length extrapolation.

The taxonomy is deliberately empirical and scoped to this dissertation. It
should not be interpreted as an exhaustive theoretical classification of all
possible Transformer failures. Its value lies in making experimental claims
more precise and in exposing which additional evidence is needed before a
particular conclusion can be drawn.

\section{Interpretation of Attention Evidence}
\label{sec:discussion-attention}

Attention statistics produced task-dependent rather than universal behaviour.
For Addition and Reverse, normalised attention entropy increased as the test
length grew. Delayed Copy, however, collapsed at the first unseen length
without a comparable entropy increase. Attention diffusion therefore cannot
serve as a general explanation of the observed failures.

The local attention ratio decreased with length across the analysed models,
but this statistic uses a fixed $\pm 10$-position neighbourhood. Part of the
decrease follows mechanically from evaluating a fixed window inside a larger
sequence. Average attention distance also increased with length, but it is
measured in raw token positions and therefore grows partly because longer
sequences permit larger distances.

These limitations illustrate why attention weights should not automatically
be treated as explanations of model predictions. Prior work has shown that
different attention distributions can sometimes produce similar outputs and
that attention weights need not correspond to other measures of feature
importance \cite{jain2019attention}. Accordingly, this dissertation treats
attention statistics as descriptive supporting evidence rather than as a
causal mechanism or a separate taxonomy category.

A stronger mechanistic study would require targeted interventions, such as
ablating particular heads, replacing learned attention patterns, or measuring
whether controlled changes in attention alter the model's output errors.

\section{Interpretation of the Real-World Extension}
\label{sec:discussion-real-world}

The SelfRegulationSCP2 experiment produced a flatter accuracy curve than the
synthetic benchmark. Across the evaluated prefix lengths and positional
encodings, mean classification accuracy remained within approximately
49--56\%. It did not show the transition from high training-length accuracy
to zero performance observed in the eligible synthetic baselines.

This contrast must be interpreted cautiously. SelfRegulationSCP2 is a
two-class EEG time-series dataset with only 200 training examples and 180 test
examples \cite{bagnall2018uea}. Held-out accuracy at the training prefix
length was already close to the 50\% uniform two-class reference. The absence
of a subsequent collapse therefore does not demonstrate successful
length-general classification.

The synthetic and real-world experiments also use different output
structures and evaluation metrics. Exact match requires every valid output
token to be correct, while time-series classification requires one class
label. In addition, longer SelfRegulationSCP2 prefixes introduce additional
observations that may be informative, irrelevant, or noisy.

The extension therefore demonstrates a flat but weak classification curve,
not evidence that real-world tasks are generally protected from
length-related failure. Its broader lesson is methodological: the validity of
the starting baseline and the meaning of the evaluation metric must be
established before comparing behaviour across lengths or task types.

\section{Limitations and Threats to Validity}
\label{sec:discussion-limitations}

\subsection{Internal Validity}

The main synthetic benchmark uses three random seeds. This is stronger than
reporting one run, but it remains a small sample for estimating a complete
performance distribution. The position-wise, attention, and mitigation
diagnostics use representative seed-42 checkpoints and therefore
characterise selected failure cases rather than multi-seed prevalence.

The mitigation experiments are not exhaustive hyperparameter searches.
Moreover, the original Delayed Copy and Reverse baselines used a different
batch size from several mitigation runs. The sinusoidal Delayed Copy
single-length control reproduced the original baseline curve, providing a
useful check for that pipeline, but equivalent controls were not run for
every task--encoding pair.

The Reverse curriculum comparison has an additional configuration mismatch:
the original baseline used a vocabulary size of 100, whereas the curriculum
configuration used 102. That experiment is therefore exploratory and cannot
be treated as a fully controlled baseline--mitigation comparison.

\subsection{Construct Validity}

Exact-match accuracy is appropriate for deterministic algorithmic tasks, but
it is deliberately strict. A single incorrect output token makes the entire
example incorrect. Token-error percentages and position-wise diagnostics
were included to prevent this strict metric from hiding partial or spatially
structured behaviour.

The first failure length is observed only on a discrete grid. The experiments
cannot determine whether collapse begins immediately above the training
length or only close to the first sampled unseen length.

Attention statistics are summary measures rather than causal tests. They are
averaged over a limited number of diagnostic batches and should not be
interpreted as full mechanistic explanations.

The positional encoding implementations also differ from some original
settings in the literature. In particular, the present Transformer is
bidirectional, and ALiBi uses a symmetric distance bias. Learned positional
entries beyond the training length exist in the model but receive no direct
training signal. Conclusions therefore apply to the implementations tested
here rather than every possible form of the named encoding.

\subsection{External Validity}

The synthetic experiments use one compact Transformer architecture and three
deterministic tasks. Larger models, decoder-only architectures, causal
attention, alternative optimisation procedures, or tasks requiring different
computations may produce different results.

Algorithmic reasoning studies have shown that Transformer length
generalisation can vary considerably with task representation, scratchpad
format, and the complexity of the required computation
\cite{anil2022exploring,zhou2024algorithms}. The present negative results
should therefore not be generalised to every Transformer or training
procedure.

The real-world extension contains only one small dataset and no specialised
time-series baseline. It provides contextual evidence rather than a broad
empirical claim about applied long-sequence classification.

\section{Implications and Future Work}
\label{sec:discussion-implications}

The principal methodological implication is that length-generalisation
evaluation should proceed in stages:

\begin{enumerate}
    \item establish that the training task is reliably learned;
    \item evaluate multiple random seeds before pooling results;
    \item report performance across a sequence of unseen lengths;
    \item examine token-level and position-wise error behaviour;
    \item use task-specific diagnostics where possible; and
    \item compare interventions with an appropriate reference baseline.
\end{enumerate}

This procedure distinguishes the ability to process a long context from the
ability to preserve a learned computation under a length shift. Architectures
and benchmarks designed for long sequences demonstrate that longer contexts
can be processed efficiently \cite{dai2019transformerxl,tay2021longrange},
but context capacity alone does not establish algorithmic length
generalisation.

Future experiments should evaluate more seeds and denser length grids,
especially between the training length and the first observed collapse.
Hyperparameters should be selected separately for baseline fitting and
mitigation evaluation, with matched batch sizes and exact configuration
checks.

The Reverse curriculum experiment should be repeated with vocabulary size 100
to provide a fully matched comparison. A direct implementation of randomised
positional indices \cite{ruoss2023randomized} would test a different
hypothesis from right-padded variable-length training. Scratchpad or
intermediate-step representations could also be investigated because their
format has been shown to affect length generalisation in reasoning tasks
\cite{anil2022exploring,kazemnejad2023impact}.

Finally, a stronger real-world extension should first establish competitive
held-out performance at the training prefix length. It should then evaluate
multiple datasets, task-specific time-series baselines, additional seeds, and
input-length interventions designed around the information structure of each
dataset.

\section{Summary}

The discussion shows that the experimental outcomes answer different
questions and should not be reduced to a single measure of long-length
accuracy. Five task--encoding configurations reliably learned their training
tasks and then collapsed at the first evaluated unseen length. Other
configurations underfit or were unstable across seeds. Position-wise
diagnostics distinguished a modest Addition error gradient from the nearly
uniform failure of Delayed Copy and Reverse.

Attention statistics provided descriptive but non-universal evidence.
Mixed-length variants underfit in the evaluated settings, while valid
curriculum and randomised-padded Delayed Copy variants produced no observed
improvement beyond the training length. The real-world extension showed a
flat but weak classification curve and therefore did not establish successful
length generalisation.

Together, these findings support an evidence-based taxonomy comprising
training-length underfitting, seed instability, immediate extrapolation
collapse, spatial error structure, and no observed improvement under tested
mitigation. The main contribution is both the observed failure behaviour and
the disciplined sequence of checks used to interpret it.

%======================================================================


\chapter{Conclusion}
\label{chap:conclusion}

\noindent
This dissertation investigated which failure behaviours occur when compact
Transformer models trained on short sequences are evaluated on substantially
longer sequences. The study used a controlled train-short/test-long framework
covering three synthetic tasks, five positional encoding settings, three
random seeds, position-wise and task-specific diagnostics, targeted
mitigation experiments, and an exploratory real-world extension.

The principal conclusion is that high training-length accuracy is not
sufficient evidence of length-general computation. All five task--encoding
configurations that satisfied the multi-seed baseline-eligibility criterion
fell to 0\% exact-match accuracy at the first evaluated unseen length.
However, the broader experiment also showed that poor long-length accuracy
does not always represent the same phenomenon. Some configurations underfit
at the training length, some were unstable across seeds, and others reliably
learned the training task before collapsing under length extrapolation.

\section{Principal Findings}

The baseline analysis identified five eligible configurations: Addition with
learned and sinusoidal positional encodings, Delayed Copy with sinusoidal
encoding and RoPE, and Reverse with learned positional embeddings. Addition
collapsed when the digit length increased from 16 to 32, while Delayed Copy
and Reverse collapsed when the sequence length increased from 128 to 256.
These outcomes were reported in Chapter~\ref{chap:synthetic-results}.

No positional encoding was uniformly successful across the three tasks.
Positional encoding affected whether the training task was learned, but
training-length success did not guarantee extrapolation. NoPE and ALiBi
frequently underfit under the evaluated architecture and optimisation
settings, while learned, sinusoidal, and rotary encodings produced
task-dependent combinations of eligibility, seed instability, and
extrapolation collapse.

The error diagnostics in Chapter~\ref{chap:failure-analysis} showed that the
form of collapse also differed across tasks. At the first unseen length,
Addition produced approximately 86.9--88.0\% token error and retained a modest
positional gradient, with lower error near the beginning and near-reference
error towards later output positions. Delayed Copy and Reverse produced
approximately 98.8--99.0\% token error that was almost uniform across the
output.

Addition errors differed by less than one percentage point between ordinary
digits with and without an incoming carry. Reverse errors were similarly high
across short and long normalised dependency distances. Therefore, Addition
collapse was not primarily concentrated on carry-dependent digits, and
Reverse collapse was not confined to the longest input--output dependencies.

The attention diagnostics produced task-dependent results. Normalised
attention entropy increased for Addition and Reverse but remained
approximately stable for Delayed Copy at the first unseen length. Attention
change was therefore retained as descriptive evidence rather than treated as
a universal explanation or an independent failure category.

The mitigation experiments also produced two distinct outcomes. Mixed-length
variants underfit in the evaluated Addition and sinusoidal Delayed Copy
settings, making their extrapolation results inconclusive. Curriculum and
randomised-padded Delayed Copy variants preserved training-length performance
but produced no observed improvement at unseen lengths. These results are
scoped negative findings about the tested configurations; they do not show
that the failures are universally resistant to mitigation.

Finally, the exploratory SelfRegulationSCP2 experiment in
Chapter~\ref{chap:real-world-extension} produced mean classification accuracy
of approximately 49--56\% across the evaluated prefix lengths. The curve was
relatively flat, but performance was already close to the uniform two-class
reference at the training prefix length. It therefore represented flat but
weak classification performance rather than successful real-world length
generalisation.

\section{Contributions}

The dissertation makes five main contributions:

\begin{enumerate}
    \item It establishes a baseline-eligibility procedure that separates
    failure to learn the training task from failure under length
    extrapolation.

    \item It provides a controlled multi-task and multi-seed comparison of
    learned, sinusoidal, NoPE, ALiBi, and RoPE positional encoding settings.

    \item It extends aggregate accuracy analysis with token-error
    percentages, position-wise profiles, carry-conditioned errors, and
    dependency-distance diagnostics.

    \item It evaluates mixed-length, curriculum, and randomised-padded
    training against corresponding reference baselines while explicitly
    checking whether each intervention preserves training-length performance.

    \item It proposes an evidence-based taxonomy that distinguishes five
    observable outcomes:
    \begin{enumerate}
        \item training-length underfitting;
        \item seed instability;
        \item immediate extrapolation collapse;
        \item spatial error structure; and
        \item no observed improvement under tested mitigation.
    \end{enumerate}
\end{enumerate}

The taxonomy is not intended as an exhaustive theory of every possible
Transformer failure. Its contribution is methodological: it provides a
structured sequence of checks that makes length-generalisation claims more
precise and prevents different experimental outcomes from being collapsed
into one statement that a model ``fails''.

\section{Limitations and Future Work}

The conclusions are limited to one compact Transformer architecture, three
synthetic tasks, five positional encoding implementations, and three baseline
seeds. Future experiments should evaluate larger and alternative
architectures, additional algorithmic tasks, more seeds, and denser test
lengths between the training length and the first observed collapse.

Mitigation experiments should use fully matched optimisation settings and
should be repeated across multiple seeds. In particular, the Reverse
curriculum experiment should be repeated using vocabulary size 100 to match
the original baseline. The effects of batch size and other optimisation
choices should also be controlled through single-length reference runs.

Future interventions could test genuinely randomised positional indices,
alternative curricula, intermediate-step or scratchpad representations, and
architectures with iterative or recurrent computation. These experiments
would test different hypotheses from the right-padded variable-length
training used in this dissertation.

A stronger mechanistic analysis should move beyond aggregate attention
statistics. Controlled head or layer ablations, attention replacement,
probing experiments, and causal interventions could test whether particular
internal changes contribute directly to output errors.

Finally, the real-world extension should be expanded only after establishing
a stronger held-out baseline at the training prefix length. Future work should
include multiple long-sequence datasets, specialised time-series baselines,
task-specific hyperparameter selection, and additional random seeds.

\section{Final Conclusion}

The answer to the primary research question is that Transformer
length-generalisation failure appears as a set of empirically distinguishable
behaviours rather than one uniform reduction in accuracy. A valid analysis
must first establish whether the training task was learned, determine whether
that result is stable across seeds, locate the first observed extrapolation
collapse, examine how errors are distributed, and test interventions against
an appropriate baseline.

Within the evaluated synthetic benchmark, every eligible configuration
collapsed at the first sampled unseen length. Yet the position-wise,
task-specific, seed-level, and mitigation evidence showed that the surrounding
failure behaviour differed substantially across configurations. The central
contribution of this dissertation is therefore not only the observation that
length extrapolation fails, but a reproducible framework for determining what
kind of experimental failure has occurred and what conclusions the available
evidence can support.





% =============================================================================

\backmatter

\bibliography{references}

% -----------------------------------------------------------------------------

\appendix

\chapter{Supplementary Position-Wise Error Heatmaps}
\label{app:position-error-heatmaps}

\noindent
Chapter~\ref{chap:failure-analysis} compared position-wise error at the
training length and the first evaluated unseen length. This appendix provides
the corresponding heatmaps across every evaluated sequence length.

The heatmaps use representative seed-42 checkpoints. Target positions are
divided into ten equal-width relative-position bins. The shared colour scale
shows token-error rate from 0\% to 100\%. These figures characterise selected
failure cases and should not be interpreted as multi-seed estimates.

\section{Addition}

\begin{figure}[p]
\centering
\includegraphics[width=\textwidth]
{addition_position_error_heatmap.pdf}
\caption{Position-wise Addition token-error rates across all evaluated digit
lengths for the learned and sinusoidal seed-42 checkpoints. Target positions
are divided into ten relative-position bins, and all panels use the same
0--100\% colour scale.}
\label{fig:appendix-addition-position-heatmap}
\end{figure}



\section{Delayed Copy}

\begin{figure}[p]
\centering
\includegraphics[width=\textwidth]
{copy_position_error_heatmap.pdf}
\caption{Position-wise Delayed Copy token-error rates across all evaluated
sequence lengths for the sinusoidal and RoPE seed-42 checkpoints. Target
positions are divided into ten relative-position bins, and all panels use the
same 0--100\% colour scale.}
\label{fig:appendix-copy-position-heatmap}
\end{figure}



\section{Reverse}

\begin{figure}[p]
\centering
\includegraphics[width=0.78\textwidth]
{reverse_position_error_heatmap.pdf}
\caption{Position-wise Reverse token-error rates across all evaluated
sequence lengths for the learned seed-42 checkpoint. Target positions are
divided into ten relative-position bins, and the colour scale shows token
error from 0\% to 100\%.}
\label{fig:appendix-reverse-position-heatmap}
\end{figure}



\section{Interpretive Scope}

The heatmaps reinforce the distinction identified in the main analysis.
Addition retains some positional variation at the first unseen length,
particularly near the beginning of the target. Delayed Copy and Reverse
instead show near-uniform high token error across positions immediately
outside the training length. At longer lengths, all three tasks approach
their respective uniform-prediction error references.

No additional statistical claims are based on these appendix figures. Their
purpose is to provide the complete position-by-length visualisation underlying
the more focused comparisons presented in the main chapter.

% =============================================================================

\end{document}
